\documentclass{article}
\usepackage[12pt]{extsizes}
\usepackage[a4paper, margin=0.5cm]{geometry}
\usepackage{indentfirst}
\usepackage[T1, T2A]{fontenc}
\usepackage[english, russian]{babel}
\usepackage{amsmath, amssymb, fleqn, cancel, graphicx, wrapfig, tabbing, color, cases}

\usepackage{amstext} % for \text macro
\usepackage{array}   % for \newcolumntype macro
\newcolumntype{L}{>{$}l<{$}} % math-mode version of "l" column type

\DeclareMathOperator{\rank}{rank}
\makeatletter
\newenvironment{sqcases}{%
  \matrix@check\sqcases\env@sqcases
}{%
  \endarray\right.%
}
\def\env@sqcases{%
  \let\@ifnextchar\new@ifnextchar
  \left\lbrack
  \def\arraystretch{1.2}%
  \array{@{}l@{\quad}l@{}}%
}
\makeatother

\DeclareRobustCommand{\divby}{\mathrel{\vbox{\baselineskip.65ex\lineskiplimit0pt\hbox{.}\hbox{.}\hbox{.}}}}

\setlength{\parindent}{0pt}

\newenvironment{rcases}
    {\left.\begin{aligned}}
    {\end{aligned}\right\rbrace}

\usepackage{hyperref}
\hypersetup{
    colorlinks,
    citecolor=black,
    filecolor=black,
    linkcolor=black,
    urlcolor=black
}




\begin{document}

\begin{titlepage}
    \title{Математический анализ}
    \author{Бабицкий Иван | М8О-111БВ-24 \\ 02.03.02 "Фундаментальная информатика и информационные технологии"}
    \date{2024/2025 учебный год, 1-й семестр}
    \maketitle
    \begin{center}
        
\includegraphics[scale=0.2]{images/mai.png}
    \end{center}
\end{titlepage}

\tableofcontents



%-------------------------------------------------------
\newpage
\paragraph{2024.09.06 (3)}
\part{Элементы теории множеств}

\section{Введение}

\subsection{Определение}
Множество --- это набор элементов.


\subsection{Определение}
Пустое множество $\varnothing$ --- это множество без элементов. Является подмножеством любого множества.


\subsection{Определение}
Универсальное множество $\mathbb{U}$ --- множество элементов рассматриваемого типа.
\begin{gather*}
    A \subset B, \text{ если } \forall x \in A \Rightarrow x \in B. \\
    \varnothing \subset A, A \subset U, A \subseteq A.
\end{gather*}


\subsection{Определение}
Логическое высказывание --- это повествовательное предложение, про которое можно сказать истинно оно или ложно.


\subsection{Определение}
Предикат --- это утверждение, зависящее от переменной, превращающееся в логическое высказывание при подстановке вместо переменной её конкретного значения.

\subsubsection{Пример}
$x > 2$ --- предикат от одной переменной.


\subsection{Определение}
Область истинности предиката --- это множество таких значений переменной, при которых этот предикат превращается в истинное высказывание.

\subsubsection{Пример}
$x \in A$, где $A$ --- область истинности этого предиката.



\section{Квантеры}

\begin{itemize}
    \item $\forall$ --- для любого
    \item $\in$ --- принадлежит
    \item $\notin, \overline{\in}$ --- не принадлежит
    \item $\exists$ --- существует
    \item $\exists!$ --- существует, и при том единственный (-ая, -ое)
    \item $\subset$ --- подмножество
    \item $\subseteq$ --- подмножество, совпадающее с множеством
    \item $\wedge$ --- и
    \item $\vee$ --- или
    \item $:$ --- таких, что  
\end{itemize}
\includegraphics[scale=0.4]{images/2024-09-06/01.png}


\section{Операции над множествами}

\subsection{Пересечение}
\[
    A \cap B = \{x: x \in A \wedge x \in B\}
\]


\subsection{Объединение}
\[
    A \cup B = \{x: x \in A \vee x \in B\}
\]


\subsection{Разность}
\[
    A \setminus B = \{x: x \in A \wedge x \notin B\}
\] 


\subsection{Симметрическая разность}
\begin{gather*}
    (A \cup B) \setminus (A \cap B) \\
    (A \setminus B) \cup (B \setminus A)
\end{gather*}


\subsection{Дополнение}
\[
    \overline{A} = \mathbb{U} \setminus A = \{x: x \notin A\}
\]



\section{Логические высказывания}

\subsection{Операции над логическими высказываниями}

\subsubsection{Отрицание}
\begin{tabular}{c|c}
    $A$ & $\overline{A}$ \\
    \hline
    $1$ & $0$ \\
    $0$ & $1$
\end{tabular}
\[
    \urcorner (A \wedge B) = \urcorner A \vee \urcorner B \text{ или } \overline{A \wedge B} = \overline{A} \vee \overline{B}
\]

\subsubsection{Конъюнкция}
\begin{tabular}{c|c|c}
    $A$ & $B$ & $A \wedge B$ \\
    \hline
    $0$ & $0$ & $0$ \\
    $0$ & $1$ & $0$ \\
    $1$ & $0$ & $0$ \\
    $1$ & $1$ & $1$ 
\end{tabular}

\paragraph{2024.09.13 (3, 4)}
\subsubsection{Дизъюнкция}
\begin{tabular}{c|c|c}
    $A$ & $B$ & $A \vee B$ \\
    \hline
    $0$ & $0$ & $0$ \\
    $0$ & $1$ & $1$ \\
    $1$ & $0$ & $1$ \\
    $1$ & $1$ & $1$ 
\end{tabular}

\subsubsection{Импликация}
\begin{tabular}{c|c|c}
    $A$ & $B$ & $A \rightarrow B$ \\
    \hline
    $0$ & $0$ & $1$ \\
    $0$ & $1$ & $1$ \\
    $1$ & $0$ & $0$ \\
    $1$ & $1$ & $1$ 
\end{tabular}

\subsubsection{Эквиваленция}
\begin{tabular}{c|c|c}
    $A$ & $B$ & $A \leftrightarrow B, A \sim B$ \\
    \hline
    $0$ & $0$ & $1$ \\
    $0$ & $1$ & $0$ \\
    $1$ & $0$ & $0$ \\
    $1$ & $1$ & $1$ 
\end{tabular}



\section{Декартово (прямое) произведение множеств}

\subsection{Определение}
Декартово (прямое) произведение множеств $A$ и $B$ --- множество упорядоченных пар, в которых 1-й элемент принадлежит множеству $A$, а 2-й элемент --- множеству $B$.
\[
    A \times B = \{(x, y): x \in A \wedge x \in B\}
\]

\subsubsection{Пример}
\begin{gather*}
    \text{Пусть } A = [1, 2], B = [2, 3] \\
    A \times B = \{(x, y): x \in [1, 2] \wedge y \in [2, 3]\}
\end{gather*}
\includegraphics[scale=0.3]{images/2024-09-13/02.png}



\section{Соответствия}

\subsection{Определение}
Соотсветствие между множествами $A$ и $B$ --- это любое подмножество декартового произведения этих элементов.


\subsection{Определение}
Пусть задано соответствие между множествами $A$ и $B$, то есть множество некоторых пар $(x, y)$. \\
Тогда образ элемента $x$ из множества $A$ --- множество таких $y$ из $B$, которые соответствуют элементу $x$.
\includegraphics[scale=0.3]{images/2024-09-13/03.png}


\subsection{Определение}
Прообраз элемента $y$ из $B$ --- множество таких $x$ из $A$, которые соответствуют элементу $y$.


\subsection{Определение}
Объединение всех прообразов --- область определения соответствий.


\subsection{Определение}
Объединение всех образов --- область значений соответствий.


\section{Элементарная классификация соответствий}

\subsection{Многомногозначное соответствие}
Существует образ и прообраз, состоящие более чем из 1 элемента. \\
\includegraphics[scale=0.3]{images/2024-09-13/04.png}


\subsection{Одномногозначное соответствие}
Образ, состоящий более чем из 1 элемента, но все прообразы состоят не более чем из 1 элемента. \\
\includegraphics[scale=0.3]{images/2024-09-13/05.png}


\subsection{Многооднозначное соответствие}
Все образы состоят не более чем из 1 элемента, но существует прообраз, состоящий более чем из 1 элемента. \\
\includegraphics[scale=0.3]{images/2024-09-13/06.png}


\subsection{Однооднозначное соответствие}
Все образы и все прообразы состоят не более чем из 1 элемента. \\
\includegraphics[scale=0.3]{images/2024-09-13/07.png}

\paragraph{Примечание}
~\\
Последние 2 типа соответствия называются функциональными соответствиями или функциями.


\subsection{Обратное соответствие}
\begin{itemize}
    \item Если $f = \{(x, y): x \in A \wedge y \in B \wedge \text{правило}\}$, то
          $f^{-1} = \{(y, x): x \in A \wedge y \in B \wedge \text{правило}\}$
    \item Если $f \subset A \times B$, то
          $f^{-1} \subset B \times A$
\end{itemize}

\subsubsection{Пример}
\[
    f: y = x^2, \text{ то } f^{-1}: x = \pm \sqrt{y}
\]

\paragraph{Примечание}
~\\
Только соответствие, обратное однозначному, также является функцией.
\\
\\
Пусть $A$ совпадает с областью определения соответствия.
\subsection{Определение}
Инъекция (однооднозначное соответствие) --- соответствие, для которого справедливо следующее высказывание:
\[
    f(x_1) = f(x_2), \text{ то } x_1 = x_2, \text{ где } x_1, x_2 \in A, f \text{ --- обозначение соответствия}
\]
\includegraphics[scale=0.3]{images/2024-09-13/08.png}


\subsection{Определение}
Сюръекция --- это функциональное соответствие, в котором область значения совпадает с множеством $B$. \\
\includegraphics[scale=0.3]{images/2024-09-13/09.png}


\subsection{Определение}
Биекция --- это отображение, являющееся одновременно инъекцией и сюръекцией (взаимно однозначное соответствие). \\
\includegraphics[scale=0.3]{images/2024-09-13/10.png}



\section{Мощность множества}

\subsection{Определение}
Равнодушные множества --- это 2 множества, между которыми можно установить взаимнооднозначное соответствие.


\subsection{Определение}
Пусть $A$ --- подмножество $B$. Будем говорить, что мощность множества $B$ больше мощности множества $A$, 
если между $A$ и $B$ нельзя установить взаимнооднозначное соответствие.


\subsection{Определение}
Счётное множество --- это множество, равномощное множеству натуральный чисел.


\subsection{Утверждение}
Множество целых чисел счётно. \\
\includegraphics[scale=0.3]{images/2024-09-13/11.png}


\subsection{Утверждение}
Объединение конечного множества чисел счётных множеств счётно \\
\begin{gather*}
    n \text{ множеств }
    \begin{cases}
        a_1, a_2, \dots, a_n \\
        b_1, b_2, \dots, b_n \\
        c_1, c_2, \dots, c_n
    \end{cases}    
\end{gather*}


\subsection{Утверждение}
Объединение счётного числа счётных множеств счётно. \\
\includegraphics[scale=0.3]{images/2024-09-13/12.png} \\
Отрезки и интервалы имеют равную мощность.


\subsection{Утверждение}
Множество точек отрезка являются более мощными, чем счётное множество.


\subsection{Определение}
Отрезок --- это множество мощностью континум.



\section{Отношения на множестве}

\subsection{Определение}
$A \times A = A^2$ --- декартов квадрат множества $A$.


\subsection{Определение}
Отношение на множестве $A$ --- соответствие, заданное на $A^2$.
\begin{itemize}
    \item $f$ --- отношение;
    \item $a f b$ --- элемент $a$ находится в отношении $f$ с элементом $b$.
\end{itemize}


\subsection{Определение}
Отношения эквивалентности --- это отношения, обладающие следующими тремя свойствами:
\begin{itemize}
    \item рефлексивность: $a f a$ --- справедливо для любого $a$;
    \item симметричность: если $a f b$, то $b f a$;
    \item транзитивность: если $a f b$ и $b f c$, то $a f c$.
\end{itemize}


\subsection{Определения (дополнительно для экзамена)}
\begin{itemize}
    \item Отношение порядка --- отношение, которое является антисимметричным и транзитивным.
    \item Отношение частичного порядка --- отношение, которое является рефлексивным, антисимметричным, транзитивным.
    \item Отношение нестрогого порядка --- отношение, которое является рефлексивным.
    \item Отношение строгого порядка --- отношение, которое является нерефлексивным и транзитивным.
\end{itemize}



%-------------------------------------------------------
\newpage
\paragraph{2024.09.20 (3)}
\part{Аксиоматика действительных чисел}

\section{Введение}

\subsection{Определение}
Действительные числа --- это множество $\mathbb{R}$, над элементами которого можно совершать операции сложение и умножение, 
между которыми установлено отношение линейного порядка и для которых выполнено свойство полноты,
подчиняющиеся следующему набору аксиом.



\section{Аксиомы}
\begin{itemize}
    \item Группа: пункты 11.1.1-11.1.3
    \item Абелева группа: пункты 11.1.1-11.1.4
\end{itemize}

\subsection{Аксиомы сложения}
\[
    a + b = c \in \mathbb{R}
\]  

\subsubsection{Существование нуля}
\[
    \exists 0 ~ \forall a \in \mathbb{R}: a + 0 = 0 + a = a
\]

\subsubsection{Существование противоположного элемента}
\[
    \forall a \in \mathbb{R} ~ \exists (-a): a + (-a) = 0
\]

\subsubsection{Ассоциативность}
\[
    (a + b) + c = a + (b + c)
\]

\subsubsection{Коммутативность}
\[
    a + b = b + a
\]


\subsection{Аксиомы умножения}
\begin{itemize}
    \item Группа: пункты 11.2.1-11.2.3
    \item Абелева группа: пункты 11.2.1-11.2.4
\end{itemize}

\subsubsection{Существование единицы}
\[
    \exists 1 ~ \forall a \in \mathbb{R}: a \cdot 1 = 1 \cdot a = a
\]

\subsubsection{Существование обратного элемента}
\[
    \forall a \in \mathbb{R}, a \neq 0 ~ \exists a^{-1} = \frac{1}{a}: a \cdot a^{-1} = a^{-1} \cdot a = 1
\]

\subsubsection{Ассоциативность}
\[
    (a \cdot b) \cdot c = a \cdot (b \cdot c)
\]

\subsubsection{Коммутативность}
\[
    a \cdot b = b \cdot a
\]


\subsection{Определение}
Поле --- это алгебраический объект, на котором введены 2 бинарные операции, подчиняющиеся 8 аксиомам выше.


\subsection{Связь сложения и умножения}

\subsubsection{Дистрибутивность}
\[
    \forall a, b, c: (a + b) \cdot c = a \cdot c + b \cdot c
\]


\subsection{Аксиомы об отношениях $\ge$}

\subsubsection{Рефлексивность}
\[
    a \ge a
\]

\subsubsection{Антисимметричность}
\[
    a \ge b \wedge b \ge a \Rightarrow a = b
\]

\subsubsection{Транзитивность}
\[
    a \ge b \wedge b \ge c \Rightarrow a \ge c
\]

\subsubsection{Определение, что данное отношение порядка является линейным порядком}
\[
    \forall a, b \in \mathbb{R}: a \ge b \vee b \ge a
\]

\paragraph{Замечание}
\begin{itemize}
    \item Если $a \ge b$, то $b \le a$.
    \item Если $a \ge b$ и $a \neq b$, то $a > b$.
\end{itemize}

\subsubsection{Связь умножения и порядка}
\[
    a \ge 0, b \ge 0 \Rightarrow a \cdot b \ge 0
\]

\subsubsection{Связь сложения и порядка}
\[
    a \ge b \Rightarrow a + c \ge b + c
\]

\subsubsection{Аксиома полноты}
Пусть $U$ и $V$ --- такие множества, что $\forall x \in U$ и $\forall y \in V$ $x \le y$. \\
Тогда $\exists c \in \mathbb{R}$, что $x \le c \le y$. 



\paragraph{2024.09.27 (3)}
\section{Ряд утверждений}

\subsection{Утверждение}
\[
    (-a) = (-1) \cdot a
\]

\subsubsection{Доказательство}
\begin{gather*}
    (-a) + a = (-1) \cdot a + a \\
    0 = (-1) \cdot a + a
\end{gather*}


\subsection{Утверждение}
\[
    a \cdot 0 = 0
\]

\subsubsection{Доказательство}
\begin{gather*}
    a(0 + 0) = a \cdot 0 + a \cdot 0 \\
    \underbrace{a \cdot 0 + (-a \cdot 0)}_{0} = a \cdot 0 + \underbrace{a \cdot 0 + (-a \cdot 0)}_{0}
\end{gather*}


\subsection{Утверждение}
\[
    a < 0, b < 0 \Rightarrow ab > 0
\]

\subsubsection{Доказательство}
\begin{gather*}
    a < 0 \\
    a + (-a) < 0 + (-a) \\
    0 < (-a) \\
    0 < (-b) \\
    (-a)(-b) > 0 \\
    (-1) \cdot a \cdot (-1) \cdot b > 0 \\
    (-1)(-1) \cdot ab > 0 \\
    1 \cdot a \cdot b > 0 \\
    0 = (-1)(-1 + 1) \\
    0 = (-1)(-1) + (-1) \cdot 1 \\
    0 = (-1)(-1) + (-1) \\
    (-1)(-1) = 1
\end{gather*}


\subsection{Утверждение}
\[
    1 > 0
\]

\subsubsection{Доказательство}
\begin{gather*}
    \text{Пусть } 1 < 0 \\
    1 < 0, 1 < 0 \Rightarrow 1 \cdot 1 > 0 \\
    1 > 0 \text{ --- противоречие }
\end{gather*}


\subsection{Утверждение}
\[
    1 \neq 0
\]

\subsubsection{Доказательство}
\begin{gather*}
    1 = 0 \\
    a \cdot 1 = a \\
    a \cdot 0 = 0
\end{gather*}


\subsection{Утверждение}
\begin{gather*}
    \text{Если } a \cdot b = 0, \text{ то или } a = 0, \text{ или } b = 0: \\
    \forall a, b \in \mathbb{R}: a \cdot b = 0 \Rightarrow a = 0 \vee b = 0
\end{gather*}

\subsubsection{Доказательство}
\begin{gather*}
    \text{Пусть } a \neq 0, b \neq 0, \text{ но } a \cdot b = 0 \\
    ab + (-a) = (-a) \\
    a(b + (-1)) = (-a) \\
    a \cdot a^{-1} (b + (-1)) = (-1) \cdot a \cdot a^{-1} \\
    b + (-1) = (-1) \\
    b = 0 \text{ --- противоречие}
\end{gather*}




%-------------------------------------------------------
\newpage
\paragraph{2024.09.27 (4)}
\part{Натуральные числа}

\section{Введение}

\subsection{Определение}
Множество, обладающее тем свойством, что наряду с элементом $a$ ему принадлежит элемент $a + 1$, называется индуктивным множеством.


\subsection{Определение}
Множеством натуральных чисел называется минимальное индуктивное множество, содержащее единицу, то есть пересечение всех индуктивных множеств, содержащих единицу.


\subsection{Определение}
Множество $A$ называется ограниченны сверху, если $\exists c \in \mathbb{R} ~ \forall x \in A: x \ge b$ \\
Такое число $b$ называется мажорантом множества $A$.


\subsection{Определение}
Множество $A$ называется ограниченным снизу, если $\exists c \in \mathbb{R} ~ \forall x \in A: x \le c$ \\
Такое число c называется минорантом множества $A$.


\subsection{Определение}
Если множество ограничено как сверху, так и снизу, его называют ограниченным множеством.


\subsection{Определение}
Число $b$ называется максимальным элементом множества $A$, если $b \in A$ и $\forall x \in A: x \le b$ 


\subsection{Определение}
Число $c$ называется минимальным элементом множества $A$, если $c \in A$ и $\forall x \in A: x \ge c$ 



\section{Существование точных граней множества}

\subsection{Определение}
Наименьшее из мажорант множества называется точной верхней гранью этого множества.


\subsection{Утверждение}
Любое ограниченное сверху множество имеет точную вернюю грань --- супремум.

\subsubsection{Доказательство}
\begin{gather*}
    \text{Пусть } A \text { --- ограниченное множество, } B \text { --- множество его мажорант, тогда} \\ 
    \forall x \in A \text{ и } \forall y \in B: x = y. \\
    \text{По аксиоме 16 } \exists c \in \mathbb{R}: \forall x \in A \text{ и } \forall y \in B: x \le c \le y, \\
    \text{ то есть } c \text{ --- мажоранта } A, \text{ но любая другая мажоранта больше } c, \text{ то есть } c = \sup A.
\end{gather*}


\subsection{Определение}
Наибольший из минорант множества называется точной нижней гранью этого множества - инфинум $c = \inf A$.


\subsection{Утверждение}
Любое ограниченное снизу множество имеет точную нижнюю грань.



\section{Принцип вложенных отрезков}

\subsection{Определение}
Последовательностью называется отображение множества натуральных чисел в множество любой природы.

\subsubsection{Пример} 
\dots

Пусть $\{I_n\}: \forall n \in \mathbb{N} ~ I_{n+1} < I_n$, тогда $\{I_n\}$ называется системой вложенных отрезков.


\subsection{Теорема | Проекция вложенных отрезков (принцип Коша Кантора)}
Пусть $\{U_n\}, n \in \mathbb{N}$ --- система вложенных отрезков, тогда $\exists c \in \mathbb{R} ~ \forall n \in \mathbb{N}: c \in U_n$. \\
Кроме того, если $U_n = [a_n ; b_n]$ и $\forall \epsilon > 0 ~ \exists n \in \mathbb{N}$, что $b_n - a_n < \epsilon$, то эта точка --- единственная.

\subsubsection{Доказательство}
\begin{gather*}
    \text{Рассмотрим множества } \{a_1, a_2, ..., a_n, ...\} \text{ и } \{b_1, b_2, ..., b_n, ...\} \\
    \text{Очевидно, что } \forall k, m: a_k < b_m \\
    \exists c ~ \forall k, m \in \mathbb{N}: a_k \le c \le b_m \Rightarrow \forall n \in \mathbb{N}: c \in U_n \\
    \text{Пусть } \forall \epsilon > 0 ~ \exists n \in \mathbb{N}: b_n - a_n < \epsilon \\
    \text{Пусть } c_1 \neq c, c_1 \in U_n; n \in \mathbb{N} \\
    \text{Обозначим } \epsilon^{-1} = |c_1 - c| \\
    \exists n \in \mathbb{N}, \text{что } b_n - a_n < \epsilon \text{, то есть } c_1 \text{ и } c \text{ не могут одновременно принадлежать отрезку } U_n
\end{gather*}

\subsubsection{Пример}
\[
U_n = \left[\frac{1}{2} - \frac{1}{n} ; \frac{1}{2} + \frac{1}{n}\right], c = \frac{1}{2}
\]


\subsection{Определение}
Система множеств $A_1$, $A_2$, ..., $A_n$, ... называестя покрытием множества $B$, если $\forall x \in B ~ \exists n: x \in A_n$.


\subsection{Определение}
Подпокрытием называется набор множеств покрытия, также являющийся покрытием.


\subsection{Пример}
\begin{gather*}
    B = \left[0, 1\right],~A_1 = \left(-\frac{1}{2} ; \frac{2}{3}\right),~A_2 = \left(\frac{1}{3} ; \frac{4}{3}\right),~A_3 = \left(\frac{1}{2} ; \frac{2}{3}\right) \\
    A_1, A_2 \text{ --- подпокрытие} \\
    A_1, A_2, A_3 \text{ --- покрытие}
\end{gather*}


\subsection{Теорема | Принцип конечного подпокрытия (принцип Бареля Лебега)}
Из любого покрытия отрезка интервалами можно выделить конечные покрытия.

\subsubsection{Пример}
\includegraphics[scale=0.4]{images/2024-09-27/01.jpg}

\subsubsection{Доказательство}
Пусть существует бесконечное покрытие отрезка, из которого нельзя выделить конечное покрытие отрезка.

Разделим пополам исходный отрезок.

Тогда для хотя бы одной половины тоже требуется бесконечность покрытия.

Разделим и его пополам и убедимся в том, что одна из новых половин также требует бесконечное покрытие.

Продолжим процесс деления пополам до бесконечности.

Заметим, что выделенная каждый раз половинка образует вложенную систему отрезков, пересечением которых будет некоторая единственная точка $C$, принадлежащая какому-то интервалу из исходного покрытия.

Так как длина половинки может быть как угодно малой, какая-то из половинок целиком попадёт внутрь этого интервала, то есть этот интервал уже является покрытием соответствующей половинки, а вовсе не бесконечного множества интервалов.


\subsection{Определение}
Окрестностью точки называется любой интервал, содержащий эту точку


\subsection{Определение}
Проколотой окрестностью точки называется окрестность точки, из которой исключена сама точка


\subsection{Определение}
Граничной точкой множества называется такая его точка, любая окрестность которой содержит как элементы множества, так и элементы, не принадлежащие множеству


\subsection{Определение}
Внутренней точкой множества называется точка множества, принадлежащая ему вместе с некоторой своей окрестностью.


\subsection{Определение}
Предельной точкой множества называется такая его точка, в любой окрестности которой содержится бесконечно много точек множества.

\subsubsection{Пример}
\includegraphics[scale=0.2]{images/2024-09-27/02.jpg}


\subsection{Определение}
Изолированной точкой множества называется такая точка этого множества, у которой существует проколотая окрестность, не содержащая элементов множеств.


\subsection{Теорема | Принцип существования}
Любое бесконечное ограниченное множество имеет предельную точку

\subsubsection{Доказательство}
$A$ --- бесконечное множество, $\exists [a, b]: A \subset [a, b]$.

Пусть $U_n$ --- произвольное покрытие отрезка $[a, b]$ интервалами, тогда $U_n$ --- покрытие множества $A$.

Пусть $A$ не имеет предельных точек, тогда любой интервал из $U_n$ содержит не более чем конечное множество точек из множества $A$.

Из $U_n$ можно выделить конечное подпокрытие, то есть $A$ --- конечное множество.




%-------------------------------------------------------
\newpage
\paragraph{2024.10.04 (3), 2024.10.11 (3, 4)}
\part{Предел последовательности}

\section{Введение}

\subsection{Определение}
Числовой последовательностью называется отображение множества $\mathbb{N}$ в множестве $\mathbb{R}$.
\[
\{a_n\}, a_n \in \mathbb{R}, n \in \mathbb{N}
\]

\subsubsection{Пример}
\begin{gather*}
    a_n = 1 \\
    a_n = (-1)^n \\
    a_n = \frac{1}{n} \\
    a_n = 1 + (-1)^n
\end{gather*}


\subsection{Определение}
Последовательность называется ограниченной, если множество её значений ограничено.


\subsection{Определение}
Число $A$ называется пределом последовательности $a_n$, если какова бы не была окрестность точки $A$, найдётся такое $N$, что $\forall n \in \mathbb{N}, n > N, a_n$ находится в этой окрестности.
\[
\lim_{n \to \infty} a_n = A
\]


\subsection{Определение}
\[
(\lim_{n \to \infty} a_n = A) \Leftrightarrow (\forall \epsilon > 0 ~ \exists N(\epsilon), \forall n \in \mathbb{N}: |a_n - A| < \epsilon)
\]


\subsection{Утверждение}
Последовательность может иметь не более одного предела

\subsubsection{Доказательство}
\begin{gather*}
    \text{Пусть } \exists \lim_{n \to \infty} a_n = A_1, \exists \lim_{n \to \infty} a_n = A_2, A_1 < A_2 \\
    \text{Пусть } \epsilon = \frac{A_2 - A_1}{3} \\
    \epsilon N_1 = \forall n > N_1: |a_1 - A_1| < \epsilon \Leftrightarrow A_1 - \epsilon < a_n < A_1 + \epsilon \\
    \epsilon N_2 = \forall n > N_2: |a_2 - A_2| < \epsilon \Leftrightarrow A_2 - \epsilon < a_n < A_2 + \epsilon \\
    \text{Пусть } N = \max\{N_1, N_2\}: \\
    \forall n > N: a_1 < A_1 + \epsilon = A_1 + \frac{A_2 - A_1}{3} = \frac{2A_1 + A_2}{3} < \frac{A_1 + 2A_2}{3} = A_2 - \epsilon < a_n    
\end{gather*}


\subsubsection{Замечание}
Ограниченная последовательность может не иметь предел

\subsubsection{Пример}
\includegraphics[scale=0.2]{images/2024-10-04/01.jpg}


\subsection{Определение}
\begin{gather*}
    (\lim_{n \to \infty} a_n = \infty) \Leftrightarrow (\forall \epsilon > 0 ~ \exists N(\epsilon) ~ \forall n > N: |a_n| > \epsilon) \\
    (\lim_{n \to \infty} a_n = +\infty) \Leftrightarrow (\forall \epsilon > 0 ~ \exists N(\epsilon) ~ \forall n > N: a_n > \epsilon) \\
    (\lim_{n \to \infty} a_n = -\infty) \Leftrightarrow (\forall \epsilon > 0 ~ \exists N(\epsilon) ~ \forall n > N: a_n > -\epsilon) \\
    (\{a_n\} \text{ является сходящей}) \Leftrightarrow (\forall \epsilon > 0 ~ \exists N(\epsilon), m \in \mathbb{N}, n > m > N: |a_n - a_m| < \epsilon)
\end{gather*}


\subsection{Определение}
Если предел последовательности равен $\infty$, $+\infty$ или $-\infty$ последовательность называется \\
бесконечно большой.


\subsection{Определение}
Последовательность называется сходящейся, если она имеет конечный (числовой) предел.


\subsection{Определение}
Если последовательность имеет предел равный нулю, то она называется бесконечно малой \\
последовательностью.


\subsection{Определение}
Последовательность $\{a_n\}$ называется фундаментальной, если 
\[
\forall\epsilon > 0, \exists{N(\epsilon)}: \forall{m} > n > N: |a_m - a_n| < \epsilon
\]


\subsection{Критерий Коши} 
(последовательность является фундаментальной) $\Leftrightarrow$ (последовательность является сходящейся)

\subsubsection{Доказательство $1 \Rightarrow 2$}
\begin{gather*}
    \exists \lim_{n\to\infty} a_n = A, A \in \mathbb{R} \\
    \forall\epsilon > 0~\exists{N(\epsilon)} \\
    \forall{n} > N, n \in \mathbb{N}: |a_n - A| < \frac{\epsilon}{2} \\
    \forall{m} > N, n \in \mathbb{N}: |a_m - A| < \frac{\epsilon}{2} \\
    |a_n - a_m| = |a_n - A - (a_m - A)| \le |a_n - A| + |a_m - A| < \frac{\epsilon}{2} + \frac{\epsilon}{2} = \epsilon \\
\end{gather*}

\subsubsection{Доказательство $1 \Leftarrow 2$}
\begin{gather*}
	\forall \epsilon > 0~\exists N(\epsilon): \forall m > n > N : |a_n - a_m| < \frac{\epsilon}{3} \\
	\text{Пусть } x_n =\inf_{m \ge n} a_m \\
	\text{Пусть } y_n =\sup_{m \ge n} a_m \\
	x_n \le x_{n+1} \le y_{n+1} \le y_n \\
	\{[x_n, y_n]\} \text{ --- система вложенных отрезков} \Rightarrow \exists c: \forall n \in \mathbb{N}, c \in [x_n, y_n] \\
	m > n \\
	a_n - \frac{\epsilon}{3} < a_m < a_n + \frac{\epsilon}{3} \\
	a_n - \frac{\epsilon}{3} \le x_m \le c \le y_m \le a_n + \frac{\epsilon}{3} \\
	|a_n - c| \le \frac{\epsilon}{3} \\
	|a_m - a_n| < \frac{\epsilon}{3} \\
	|x_m - a_n| \le \frac{\epsilon}{3} \\
	|y_m - a_n| \le \frac{\epsilon}{3} \\
	|a_m - c| \le |a_m - a_n + a_n - c| \le |a_m - a_n| + |a_n - c| \le
	\frac{\epsilon}{3} + \frac{\epsilon}{3} = \frac{2\epsilon}{3} < \epsilon
\end{gather*}

\includegraphics[scale=0.2]{images/2024-10-04/02.jpg}

\subsubsection{Следствие}
Область значений последовательности есть ограниченное множество, то есть последовательность ограничена.



\section{Алгебраические свойства пределов последовательности}
Пусть $\lim_{n \to \infty} a_n = A, \lim_{n \to \infty} b_n = B$

\subsection{Утверждение | Сложение}
$\exists \lim_{n \to \infty} (a_n + b_n) = A + B$

\subsubsection{Доказательство}
\begin{gather*}
	\text{Если } (\lim_{n \to \infty} a_n = A) \Leftrightarrow (\forall \epsilon > 0 \text{ } \exists N_1(\epsilon): \forall n > N_1, n \in \mathbb{N} : |a_n - A| < \frac{\epsilon}{2}) \\
	(\lim_{n \to \infty} b_n = B) \Leftrightarrow (\forall \epsilon > 0 \text{ } \exists N_2(\epsilon): \forall n > N_2, n \in \mathbb{N} : |b_n - B| < \frac{\epsilon}{2}) \\
	\text{Пусть } N = \max\{N_1, N_2\} \\
	\forall n > N : |a_n + b_n - (A + B)| \le |a_n - A| + |b_n - B| \le \frac{\epsilon}{2} + \frac{\epsilon}{2} = \epsilon \Leftrightarrow \lim_{n \to \infty} (a_n + b_n) = A + B
\end{gather*}


\subsection{Утверждение | Умножение}
$\exists \lim_{n \to \infty} (a_n \cdot b_n) = A \cdot B$

\subsubsection{Доказательство}
\begin{gather*}
	a_n \cdot b_n = (a_n - A + A)(b_n - B + B) = (a_n - A)(b_n - B) + A(b_n - B) + B(a_n - A) + AB \\
	\forall \epsilon > 0 ~ \exists N_1(\epsilon): \forall n > N_1 : |a_n - A| < \epsilon \\
	\forall \epsilon > 0 ~ \exists N_2(\epsilon): \forall n > N_2 : |b_n - B| < \epsilon \\
	\forall \epsilon > 0 ~ \exists N_3(\epsilon): \forall n > N_3 : |b_n - B| < \frac{\epsilon}{3|A|} \text{ (если $A \neq 0$, иначе отсутствует)} \\
	\forall \epsilon > 0 ~ \exists N_4(\epsilon): \forall n > N_4 : |a_n - A| < \frac{\epsilon}{3|B|} \text{ (если $A \neq 0$, иначе отсутствует)} \\
	\text{let } N = \max\{N_1, N_2, N_3, N_4\} \\
	\forall n > N: |a_n b_n - AB| \le |a_n - A||b_n - B| + |A||b_n - B| + |B||a_n - A| < \\
	< \epsilon^2 + |A| \frac{\epsilon}{3|A|} + |B| \frac{\epsilon}{3|B|} < 3 \cdot \frac{\epsilon}{3} = \epsilon \\
	(\epsilon^2 < \frac{\epsilon}{3} \Leftarrow \epsilon < \frac{1}{3})
\end{gather*}

\subsection{Утверждение | Деление}
Если $B \neq 0$, то $\exists \lim_{n \to \infty} \frac{a_n}{b_n} = \frac{A}{B}$

\subsubsection{Доказательство}
\begin{gather*}
\frac{a_n}{b_n} = \frac{A}{B} = \frac{a_n \cdot B - b_n \cdot A}{b_n \cdot B} = \frac{(a_n - A) \cdot B + A \cdot B - (b_n - B) \cdot A - A \cdot B}{b_n \cdot B} = \frac{a_n - A}{b_n} - \frac{b_n - B}{b_n} \cdot \frac{A}{B} \\
\exists \epsilon ~ N_1 \forall n > N_1 |b_n| > \frac{B}{2} \\
\exists \epsilon ~ N_2 \forall n > N_2 |a_n - A| < \frac{\epsilon|A|}{4} \\
\exists \epsilon ~ N_3 \forall n > N_3 |b_n - B| < \frac{\epsilon|B|^2}{4|A|} \\
\text{Пусть } N = max\{N_1, N_2, N_3\}, \forall n > N: \\
\frac{|a_n - A|}{b_n - B} \le {} \frac{|a_n - A|}{|b_n|} + \frac{|b_n - B|}{b_n} \cdot \frac{|A|}{|B|} \le \frac{2 \cdot \epsilon \cdot |B|}{4|B|} + \frac{2 \cdot \epsilon* |B|^2}{4|B|} \cdot \frac{|A|}{|B|} = \frac{\epsilon}{2} + \frac{\epsilon}{2} = \epsilon
\end{gather*}



\section{Предельный переход в неравенства}

\subsection{Утверждение}
\begin{gather*}
    \text{1. } \exists \lim_{n \to \infty} a_n = A \\
    \text{2. } \exists \lim_{n \to \infty} b_n = B \\
    \text{3. } \forall n \in \mathbb{N}: a_n \le c_n \le b_n
\end{gather*}

\subsubsection{Доказательство $\Rightarrow$}
$\exists \lim_{n \to \infty} c_n = A$



\subsection{Утверждение}
\begin{gather*}
    \text{1. } \exists \lim_{n \to \infty} a_n = A \\
    \text{2. } \exists \lim_{n \to \infty} b_n = B \\
    \text{3. } \forall n \in \mathbb{N}: a_n < b_n
\end{gather*}

\subsubsection{Доказательство $\Rightarrow$}
\begin{gather*}
    \forall \epsilon > 0 ~ \exists N_1(\epsilon): \\
    \forall n > N_1: |a_n - A| < \epsilon \Leftrightarrow A - \epsilon < a_n < A + \epsilon \\
    \forall n > N_2: |b_n - B| < \epsilon \Leftrightarrow B - \epsilon < b_n < B + \epsilon \\
    c_n = b_n - a_n, \exists \lim_{n \to \infty} c_n = B - a, c_n \ge 0, \forall n \in \mathbb{N} \\
    \text{Пусть } B - A < 0: \epsilon = \frac{1}{2} \cdot |B - A| \\
    \exists N ~ \forall n > N: c_n < B - A + \epsilon = B - A + \frac{1}{2} \cdot (A - B) = \frac{1}{2} \cdot B - \frac{1}{2} \cdot A < 0 \\
    b_n - a_n < 0 \text{ --- противоречие}
\end{gather*}


\subsection{Утверждение | Принцип двух полицейских}
\begin{gather*}
	\begin{cases}
		\exists \lim_{n \to \infty} a_n = A \\
		\exists \lim_{n \to \infty} b_n = A \\
		\forall n \in \mathbb{N}: a_n \le c_n \le b_n
	\end{cases} \Rightarrow \lim_{c \to \infty} c_n = A
\end{gather*}

\subsubsection{Доказательство}
\begin{gather*}
    \forall \epsilon ~ \exists N(\epsilon) ~ \forall n > N_1 \\
    \forall \epsilon ~ \exists N(\epsilon) ~ \forall n > N_2 \\
    N = \max{N_1, N_2} \\
    \forall n > N: A - \epsilon < a_n \le c_n \le b_n < A + \epsilon \\
    A - \epsilon < c_n < A + \epsilon \\
    |c_n - A| < \epsilon
\end{gather*}



\section{Монотонные последовательности}

\subsection{Определение}
Последовательность $\{a_n\}$ называется возрастающей, если для $\forall n \in \mathbb{N} ~ a_n < a_n+1$.


\subsection{Определение}
Последовательность $\{a_n\}$ называется неубывающей, если для $\forall n \in \mathbb{N} ~ a_n \le a_n+1$.


\subsection{Определение}
Последовательность $\{a_n\}$ называется убывающей, если для $\forall n \in \mathbb{N} ~ a_n > a_n+1$.


\subsection{Определение}
Последовательность $\{a_n\}$ называется невозрастающей, если для $\forall n \in \mathbb{N} ~ a_n \ge a_n+1$.


\subsection{Теорема | Вейерштрасс}
Монотонная ограниченная последовательность имеет конечный предел.

\subsubsection{Доказательство}
\begin{gather*}
    \text{Пусть для последовательности } \{a_n\}: \\
    \forall{n} \in \mathbb{N}: a_{n+1} \ge a_n, \\
    \forall{n} \in \mathbb{N}: a_n < A, \text{тогда} \\
    \text{1. } a_n \text{ --- финальная постоянная последовательность, то есть } \exists{N} ~ \forall{n} > N: a_n = a \in \mathbb{R} \\
    \lim_{n\to\infty} a_n \\
    \text{2. Множество значений } a_n \text{ бесконечно, то есть множество значений имеет предельную точку,} \\ 
    \text{так как $a_n$ --- ограничена} \\
    \text{3. } \exists c = \sup_{n \in \mathbb{N}} a_n ~ \forall \epsilon > 0 ~ \exists a_n: c - \epsilon < a_n \le c, \forall m > n: \\
    c - \epsilon < a_n \le a_m \le c \rightarrow \lim_{n \to \infty} a_n = c
\end{gather*}


\subsection{Утверждение}
Неубывающая последовательность, имеющая предел, ограничена сверху.


\subsection{Утверждение}
Для того, чтобы монотонная последовательность имела предел, необходимо и достаточно, чтобы она была ограничена.

\subsubsection{Пример}
\begin{gather*}
a_n = \frac{n}{q^n}, q > 1 \\
\frac{a_{n+1}}{a_n} = \frac{n + 1}{q^{n + 1}} \cdot \frac{q^n}{n} = \frac{n + 1}{q^n} < 1 \\
n + 1 < q \cdot n \\
1 < (q - 1) \cdot n \\
n > \frac{1}{q - 1} \text{ --- начиная с этого } n \text{, последовательность убывает} \\
a_n \text{ --- финально убывающая ограниченная снизу последовательность}
\end{gather*}

\paragraph{2024.10.18 (3)}
\subsubsection{Пример}
\begin{gather*}
    a_n = (1 + \frac{1}{n})^n \\
    (1 + \alpha)^n \ge 1 + \alpha n, \alpha > -1 \\
    \text{Пусть } n = 1, \text{ тогда } 1 + \alpha \ge 1 + \alpha \\
    \text{Пусть } n = k, \text{ тогда } (1 + \alpha)^k \ge 1 + \alpha k \\
    \text{Пусть } n = k + 1, \text{ тогда } (1 + \alpha)^{k + 1} = (1 + \alpha)^k (1 + \alpha) \ge (1 + \alpha k)(1 + \alpha) = 1 + \alpha (k + 1) + \alpha^2k > 1 + \alpha (k + 1) \\
\end{gather*}

\subsubsection{Пример}
\begin{gather*}
    b_n = \left(1 + \frac{1}{n}\right)^{n + 1} \\
    \frac{b_{n-1}}{b_n} = 
    \frac{\left(1 + \frac{1}{n - 1}\right)^n}{\left(1 + \frac{1}{n}\right)^{n + 1}} = 
    \left(\frac{1 + \frac{1}{n - 1}}{1 + \frac{1}{n}}\right)^n \frac{1}{1 + \frac{1}{n}} = 
    \left(\frac{n^2}{n^2 - 1}\right)^n \frac{n}{n + 1} = 
    \left(1 + \frac{1}{n^2 - 1}\right)^n \frac{n}{n + 1} \ge 
    \left(1 + \frac{n}{n^2 - 1}\right) \frac{n}{n + 1} > \\
    > \left(1 + \frac{1}{n}\right) \left(1 - \frac{1}{n + 1}\right) =
    1 + \frac{1}{n} - \frac{1}{n + 1} - \frac{1}{n(n + 1)} =
    1 + \frac{1}{n(n + 1)} - \frac{1}{n(n + 1)} = 1 \\
    \frac{n}{n^2 - 1} > \frac{1}{n} \\
    \frac{b_{n - 1}}{b_n} > 1 \Rightarrow b_n < b_{n - 1} \Rightarrow b_n \text{ --- убывающая}
\end{gather*}

\section{Подпоследовательности}

\subsection{Определение}
Любая бесконечная выборка из последовательности называется подпоследовательностью.
\[
    \{a_n\} \rightarrow a_{n_1}, a_{n_2}, \dots, a_{n_k}
\]

\subsubsection{Пример}
\begin{gather*}
    a_n = (-1)^n: \\
    a_2, a_4, a_6,\dots, a_{2k}, \dots \\
    a_1, a_3, a_5,\dots, a_{2k - 1}, \dots
\end{gather*}

\subsection{Утверждение}
Из любой ограниченной последовательности можно выбрать сходящуюся последовательность.

\subsubsection{Доказательство}
\begin{enumerate}
    \item Пусть область значений последовательности $a_n$ --- конечное множество. \\
          Выберем то из значений последовательности, которому соответствует бесконечно много элементов последовательности. \\
          Эти элементы и образуют исходную подпоследовательность.
    \item Если множество значений последовательности $a_n$ бесконечно, то оно имеет предельную точку. \\
          Пусть $c$ --- предельная точка.

          $\epsilon = 1$, тогда $\exists a_{k_1} \in (c - 1 ; c + 1)$

          $\epsilon = \frac{1}{2}$, тогда $\exists a_{k_2} \in (c - \frac{1}{2} ; c + \frac{1}{2}), k_2 > k_1$

          $\epsilon = \frac{1}{3}$, тогда $\exists a_{k_3} \in (c - \frac{1}{3} ; c + \frac{1}{3}), k_3 > k_2$

          ${a_{k_n}}: \lim_{n \to \infty} a_{k_n} = C$ \\
\end{enumerate}



\section{Нижний и верхний пределы последовательности}

\subsection{Определение}
Предел сходящейся подпоследовательности называется частичным пределом последовательности.

\subsubsection{Пример}
\begin{gather*}
    \text{Пусть } a_n \text{ --- ограниченная последовательность} \\
    \text{Пусть } x_n = \inf_{m \ge n} a_m, y_n = \sup_{m \ge n} a_m \\
    x_{n+1} \ge x_n, \{x_n\} \text{ --- неубывающая последовательность} \\
    \{x_n\} \text{ ограничена сверху (любое $y_k$ является мажорантом)} \\
    x_n \text{ --- сходится, } \exists \lim_{n \to \infty} x_n = m
\end{gather*}

\subsection{Определение | Нижний предел}
Нижним пределом последовательности $\{a_n\}$ называестя
\[
\exists \varliminf_{n \to \infty} a_n = \lim_{n \to \infty} x_n
\]


\subsection{Определение | Верхний предел}
Верхним пределом последовательности $\{a_n\}$ называется
\[
\exists \varlimsup_{n \to \infty} a_n = \lim_{n \to \infty} y_n
\]


\subsection{Определение}
Если последовательность не ограничена сверху, будем говорить, что
\[
\varlimsup_{n \to \infty} a_n = +\infty
\]


\subsection{Определение}
Если последовательность не ограничена снизу, будем говорить, что
\[
\varliminf_{n \to \infty} a_n = -\infty
\]


\subsection{Утверждение}
Пусть последовательность $\{a_n\}$ сходится. В этом случае верхний и нижний предел совпадают и равны просто пределу последовательности.

\subsubsection{Доказательство}
\begin{gather*}
    \text{Пусть } \exists \lim_{n \to \infty} a_n = A \Rightarrow \varlimsup_{n \to \infty} a_n = \varliminf_{n \to \infty} A\\
    x_n = \inf_{m \ge n} a_m \\
    \exists \lim x_n = A_1, A - A_1 < \epsilon, \epsilon > 0, \epsilon \in \mathbb{R} \\
    A_1 = A
\end{gather*}


\subsection{Утверждение}
Если верхний и нижний пределы последовательности равны между собой, то последовательность сходится именно к этому пределу.


\subsection{Утверждение}
Нижний предел последовательности совпадает с наименьшим из частичных пределов.


\subsection{Утверждение}
Верхний предел последовательности совпадает с наибольшим из частичных пределов.




%--------------------------------------------------
\newpage
\paragraph{2024.10.25 (3, 4)}
\part{Предел функции}

\section{Введение}

\subsection{Определение}
Базой называется набор множеств, обладающий следующими свойствами

\begin{enumerate}
    \item Пустое множество не является элементом этого набора: $\varnothing \notin \mathcal{B}$
    \item Если $B_1 \in \mathcal{B}$, $B_2 \in \mathcal{B}$, то $\exists B_3 \in \mathcal{B}: B_3 \subseteq B_1 \cap B_2$ 
\end{enumerate}

\subsubsection{Пример 1 | База проколотых окрестностей точки $x_0$}
Все интервалы, содержащие точку $x_0$, из которых исключена сама точка $x_0$. \\
\includegraphics[scale=0.3]{images/2024-10-25/01.jpg}

\subsubsection{Пример 2 | База окрестностей бесконечности}
\[
    B = \mathbb{R} \setminus \left[ a, b \right], a <0, b > 0
\]
\includegraphics[scale=0.3]{images/2024-10-25/02.jpg}

\subsubsection{Пример 3 | База окрестностей $+\infty$}
\[
    B = (a, +\infty)
\]
\includegraphics[scale=0.3]{images/2024-10-25/03.jpg}

\subsubsection{Пример 4 | $x \rightarrow x_0+$}
\[
    B = (x_0, x_0 + a), a > 0
\]
\includegraphics[scale=0.3]{images/2024-10-25/04.jpg}

\subsection{Определение}
Пусть $V(A)$ --- окрестность точки $A$.
\[
    \left( \lim_B f(x) = A \right) \overset{def}{\Leftrightarrow} \left( \forall V(A)~\exists B \in \mathcal{B}: f(B) \subset V(A) \right)
\]

\subsubsection{Перевод на другой математический язык}
\[
    B: x \to a
\] 
(фото 2)


\subsection{Определение (классическое определение предела функции)}
\[
    \left( \lim_{x \to a} f(x) = A \right) \overset{def}{\Leftrightarrow} \left( \forall \epsilon > 0~\exists \delta(\epsilon) > 0~\forall x, x \neq a, |x - a| < \delta: |f(x) - A| < \epsilon \right)
\]
(фото 3)

\subsubsection{Пример}
\begin{gather*}
    \text{Дано: } \\ 
    \lim_{x \to 2} x^2 = 4, \forall \epsilon > 0 \\
    \text{Доказать: } \\ 
    \exists \delta: \forall x \neq 2, |x - 2| < \delta (2 - \delta < x < 2 + \delta): |x^2 - 4| < \epsilon \\
    \text{Доказательство:} \\
    |x^2 - 4| < \epsilon \Leftrightarrow 4 - \epsilon < x^2 < 4 + \epsilon \\
    \sqrt{4 - \epsilon} < x < \sqrt{4 + \epsilon} ~~~~~~~~~~~~~~~~~~~~ \delta < \frac{\epsilon}{5}\\ 
    \sqrt{4 - \epsilon} < 2 - \delta, \sqrt{4 - \epsilon} > 2 + \epsilon \\
    4 - \epsilon < 4 - 4 \delta + \delta^2 \Rightarrow 4\delta - \delta^2 < \epsilon, 4 \delta + \delta^2 < \epsilon \rightarrow \frac{4}{5} \epsilon + \frac{\epsilon^2}{25} < \epsilon \Leftrightarrow \frac{\epsilon}{5} < 1 \Rightarrow \epsilon < 5 \\
    \forall \epsilon > 0~\exists \delta(\epsilon) = \min \left( \frac{\epsilon}{5} ; 1 \right): |x - 2| < \delta \Rightarrow |x^2 - 4| < \epsilon
\end{gather*}

\subsubsection{Пример}
\begin{gather*}
    y = \text{sign} x = 
    \begin{cases}
        1, x > 0 \\
        0, x = 0 \\
        -1, x < 0
    \end{cases}
    \text{(фото 4)}
\end{gather*}


\subsection{Определение}
\[
(\exists \lim_{x \to a+} f(x) = A) \overset{def}{\Leftrightarrow} (\forall \epsilon >0~\exists \delta(\epsilon) > 0~\forall x: a < x < a + \delta: |f(x) - A| < \epsilon)
\]
(фото 5)


\subsection{Определение}
Функция $f(x)$ называется финально постоянной при базе $\mathcal{B}$, если
\[
\exists B \in \mathcal{B}~\forall x \in B: f(x) = C = \text{const}, c \in \mathbb{R}
\]

\subsection{Утверждение}
\[
\lim_{\mathcal{B}} f(x) = C
\]

\subsection{Определение}
Функция $f(x)$ называеся финально ограниченной при базе $\mathcal{B}$, если
\[
\exists B \in \mathcal{B}, f(x) \text{ ограничена на } \mathcal{B}, \text{ то есть } \exists M \in \mathbb{R}, \forall x \in B.~|f(x)| < M
\]

\subsection{Определение}
Функция $f(x)$ называеися бесконечно-малой при базе $\mathcal{B}$, если
\[
\exists \lim_{\mathcal{B}} f(x) = 0
\]

\subsection{Утверждение}
\begin{gather*}
    \text{Пусть } f(x) \text{ --- бесконечно малая при } \mathcal{B} \\
    ~~~~~ g(x) \text{ --- финально-ограниченная при } \mathcal{B} \\
    \text{Тогда } f(x) \cdot g(x) \text{ --- бесконечно мало при } \mathcal{B} \\
\end{gather*}

\subsubsection{Доказательство}
\begin{gather*}
    \exists B_1 \in \mathcal{B} \text{ и } \exists M \in \mathbb{R}: \forall x \in B_1: |g(x)| < M \\
    \forall \epsilon > 0~\exists B_2 \in \mathcal{B}~\forall x \in B_2: |f(x)| < \frac{\epsilon}{M} \\
    \exists B_3 \in \mathcal{B}, B_3 \subset B_1 \cap B_2 \Rightarrow \forall x \in B_3: |f(x) \cdot g(x)| = |f(x)| \cdot |g(x)| < M \cdot \frac{\epsilon}{M} = \epsilon, \\
    \text{то есть } \lim_{\mathcal{B}} f(x) \cdot g(x) = 0
\end{gather*}

\subsubsection{Пример}
(график y = sin(1/x))


\subsection{Утверждение}
Функция при выбранной базе может иметь не более одного предела.

\subsubsection{Доказательство}
\begin{gather*}
    \text{Пусть } \exists \lim_{\mathcal{B}} f(x) = A, \exists \lim_{\mathcal{B}} f(x) = B, A \neq B \\
    V_1(A) \text{ --- окрестность } A, V_2(A) \text{ --- окрестность } B, V_1 \cap V_2 = \varnothing \\
    \exists B_1 \in \mathcal{B}, f(B_1) \subset V_1(A) \\
    \exists B_2 \in \mathcal{B}, f(B_2) \subset V_2(B) \\
    B_3 \in \mathcal{B}, B_3 \in B_1 \cap B_2 \\
    f(B_3) \subset V_1(A), f(B_3) \subset V_2(B) \text{ --- chooshe}
\end{gather*}


\subsubsection{Пример}
\begin{gather*}
    ...    
\end{gather*}


\subsection{Определение | Альтернативное классическое определение предела (Гейне)}
\[
    (\lim_{x \to a} f(x) = A) \overset{def}{\Leftrightarrow} (\forall \{x_n\}, \lim_{n \to \infty} x_n = a: \lim_{n \to \infty} f(x_n) = A)
\]

\subsection{Утверждение | Эквивалентность определений Коши и Гейне}

\subsubsection{Доказательство $\Rightarrow$}
\begin{gather*}
    \forall \epsilon > 0~\exists \delta(\epsilon) > 0~\forall x \neq a, |x - a| < \delta: |f(x) - A| < \epsilon \\
    \forall x_n \lim_{n \to \infty} x_n = a~~~~~\forall \delta > 0~\exists N(\delta)~\forall n > N: |x_n - a| < \delta \\
    \forall n > N: |f(x_n) - A| < \epsilon \Rightarrow \lim_{n \to \infty} f(x_n) = A
\end{gather*}
% 

\subsubsection{Доказательство $\Leftarrow$}
(фото)

\subsection{Утверждение}
Финально постоянная функция (при $\mathcal{B}$) имеет предел по этой базе.

\subsubsection{Доказательство}
$\exists B \in \mathcal{B}~\forall x \in B: f(x) = A, lim_{\mathcal{B}} f(x) = A$

\subsection{Утверждение}
\[
\text{Если } \lim_{\mathcal{B}} f(x) = A, \text{ то } f(x) = A \text{ --- бесконечно мала при } \mathcal{B}.
\]


\section{Арифметические свойства пределов функции}
\[
    \lim_{\mathcal{B}} f(x) = A, \lim_{\mathcal{B}} g(x) = B
\]

\subsection{Утверждение | Сложение}
\[
    \exists \lim_{\mathcal{B}} (f(x) + g(x)) = A + B
\]

\subsubsection{Доказательство}
\begin{gather*}
    \forall \epsilon > 0 \\
    \exists B_1 \in \mathcal{B}: \forall x \in B_1: |f(x) - A| < \frac{\epsilon}{2} \\
    \exists B_2 \in \mathcal{B}: \forall x \in B_2: |g(x) - B| < \frac{\epsilon}{2} \\
    \dots
\end{gather*}


\subsection{Утверждение | Умножение}
\[
    \exists \lim_{\mathcal{B}} (f(x) \cdot g(x)) = A \cdot B
\]

\subsubsection{Доказательство}
Докажем, что $f(x) \cdot g(x) - A \cdot B$ --- бесконечно мала при $\mathcal{B}$
\begin{gather*}
    f(x) \cdot g(x) - A \cdot B = 
    \underbrace{\underbrace{\underbrace{(f(x) - A)}_{\text{б.м.}} \underbrace{(g(x) - B)}_{\text{б.м.}}}_{\text{б.м.}} + 
    \underbrace{\underbrace{B}_{\text{огр.}} \underbrace{(f(x) - A)}_{\text{б.м.}}}_{\text{б.м.}} + 
    \underbrace{\underbrace{A}_{\text{огр.}} \underbrace{(g(x) - B)}_{\text{б.м.}}}_{\text{б.м.}}}_{\text{б.м.}}
\end{gather*}


\subsection{Утверждение | Умножение}
Если $B \neq 0$, то 
\[
    \exists \lim_{\mathcal{B}} \frac{f(x)}{g(x)} = \frac{A}{B}
\]

\subsubsection{Доказательство}
\begin{gather*}
    \text{Если } B \neq 0, \text{ то} \\
    \frac{f(x)}{g(x)} - \frac{A}{B} = 
    \frac{f(x) \cdot B - g(x) \cdot A}{B \cdot g(x)} = \\
    = \frac{(f(x) - A) \cdot B - (g(x) - B) \cdot A}{B \cdot g(x)} =
    \underbrace{(f(x) - A)}_{\text{б.м.}} \underbrace{\left( \frac{1}{g(x)} \right)}_{\text{б.м.}} - 
    \underbrace{(g(x) - B)}_{\text{б.м.}} \underbrace{\left( \frac{A}{B \cdot g(x)} \right)}_{\text{б.м.}} \text{ --- б.м.}
\end{gather*}
\[
\exists B_1 \in \mathcal{B}~\forall x \in B_1: \frac{|B|}{2} < |g(x)| < \frac{3|B|}{2} \Rightarrow \frac{1}{g(x)} \text{ --- финально ???}
\]


\paragraph{2024.11.01 (3)}
\subsection{Утверждение}
\begin{gather*}
    \text{Пусть } \exists \lim_{\mathcal{B}} f(x) = a, \exists \lim_{\mathcal{B}} f(x) = b, b > a \\
    \text{Тогда } \exists B_1 \in \mathcal{B} ~ \forall x \in B_1: f(x) < g(x)
\end{gather*}

\subsubsection{Доказательство}
\begin{gather*}
    \text{Пусть } \epsilon = \frac{b - a}{3} \\
    \text{Тогда }
    \begin{rcases} 
        & \exists B_2 \in \mathcal{B} ~ \forall x \in B_2: |f(x) - a| < \epsilon \Rightarrow f(x) < a + \epsilon \\
        & \exists B_3 \in \mathcal{B} ~ \forall x \in B_3: |g(x) - b| < \epsilon \Rightarrow b - \epsilon < g(x)
    \end{rcases}
    a + \epsilon < b - \epsilon \\
    \text{Пусть } B_1 \subset B_2 \cap B_3, B_1 \in \mathcal{B} \Rightarrow \forall x \in B_1: f(x) < a + \epsilon < b - \epsilon < g(x)
\end{gather*}


\subsection{Утверждение}
\begin{gather*}
    \text{Пусть } \exists \lim_{\mathcal{B}} f(x) = a, \exists \lim_{\mathcal{B}} f(x) = b, b > a \\
    \exists B_1 \in \mathcal{B} ~ \forall x \in B_1: f(x) \le g(x) \\
    \text{Тогда } a \le b
\end{gather*}

\subsubsection{Доказательство}
\begin{gather*}
    \text{Пусть } a > b, \epsilon = \frac{a - b}{3} \\
    \text{Тогда } \\
    \exists B_2 \in \mathcal{B} ~ \forall x \in B_2: |f(x) - a| < \epsilon \\
    \exists B_3 \in \mathcal{B} ~ \forall x \in B_3: |g(x) - b| < \epsilon \\
    \text{Пусть } B_4 \subset B_2 \cap B_3, B_4 \in \mathcal{B} \Rightarrow \forall x \in B_4: g(x) < b + \epsilon < a - \epsilon < f(x) \Rightarrow g(x) < f(x) \\
    \text{Пусть } B_5 \subset B_1 \cap B_4 \Rightarrow \forall x \in B_5: f(x) \le g(x), g(x) < f(x) \text{ --- reductio ad absurdum}
\end{gather*}

\subsubsection{Пример}
\begin{gather*}
    f(x) = -\sin^2 x \\
    g(x) =  \sin^2 x \\
    \text{Если } x \in [-1; 1] \setminus \{0\} ~~~ f(x) < g(x) \\
    \lim_{x \to 0} f(x) = 0 \\
    \lim_{x \to 0} g(x) = 0
\end{gather*}


\subsection{Утверждение}
\begin{gather*}
    \text{Пусть } \lim_{\mathcal{B}} f(x) = a, \lim_{\mathcal{B}} g(x) = a \\
    \exists B_1 \in \mathcal{B} ~ \forall x \in B_1: f(x) \le h(x) \le g(x) \\
    \text{Тогда } \exists \lim_{\mathcal{B}} h(x) = a
\end{gather*}

\subsubsection{Доказательство}
\begin{gather*}
    \exists B_2 \in \mathcal{B} ~ \forall x \in B_2: a - \epsilon < f(x) \\
    \exists B_3 \in \mathcal{B} ~ \forall x \in B_3: g(x) < a + \epsilon \\
    \text{Пусть } B_4 \subset B_1 \cap B_2 \cap B_3, B_4 \in \mathcal{B} ~ \forall x \in B_4: a - \epsilon < f(x) \le h(x) \le g(x) < a + \epsilon \\
    a - \epsilon < h(x) < a + \epsilon \\
    |h(x) - a| < \epsilon \\
    \lim_{\mathcal{B}} h(x) = a
\end{gather*}



\section{Первый замечательный предел}

\subsection{Определение}
\[
    \lim_{x \to 0} \frac{\sin x}{x} = 1
\]
\includegraphics[scale=0.1]{images/2024-11-01/01.jpg}
\begin{gather*}
    S_{\Delta OAB} < S_{\Delta OAC} < S_{\Delta ODC} \\
    \frac{1}{2} \cdot \cos x \cdot \sin x < \frac{x}{2} < \frac{1}{2} \cdot \tg x \\
    \cos x \cdot \sin x < x < \tg x \\
    \underbrace{cos x}_{1 \text{ при } x \rightarrow 0+} < \frac{x}{\sin x} < \underbrace{\frac{1}{\cos x}}_{1 \text{ при } x \rightarrow 0+} \\
    \lim_{x \to 0+} \frac{x}{\sin x} = 1 \Rightarrow \lim_{x \to 0} \frac{\sin x}{x} = 1 
\end{gather*}

\subsubsection{Пример}
\begin{gather*}
    y = f(x) = sin \frac{1}{x} \\
    \nexists \lim_{x \to 0} f(x)
\end{gather*}


\subsection{Определение}
\begin{gather*}
    \text{Пусть } f(x): X \rightarrow \mathbb{R} \\
    \text{Колебанием функции } f(x) \text{ на множестве } X \text{ называется } \omega(t, x) = \sup_X |f(x_1) - f(x_2)|, \forall x_1, x_2 \in X
\end{gather*}

\subsubsection{Замечание}
\[
    \text{Если } f(x) \text{ неограничена на } X \text{, будем считать, что } \omega(f, x) = +\infty
\]



\section{Критерий Коши существования предела функции}

\subsection{Утверждение}
\begin{gather*}
    \left( \exists \lim_{\mathcal{B}} f(x) = a, a \in \mathbb{R} \right) \Leftrightarrow 
    \left( \forall \epsilon > 0 ~ \exists B \in \mathcal{B}, \omega(f, B) < \epsilon \right)
\end{gather*}

\subsubsection{Доказательство 1 $\Rightarrow$ 2}
\begin{gather*}
    \exists \lim \Rightarrow \text{Критерий Коши} \\
    \forall \epsilon > 0 ~ \exists B \in \mathcal{B} ~ \forall x \in B: |f(x) - a| < \frac{\epsilon}{3} \\
    \text{Пусть } x_1, x_2 \in B \\
    |f(x_1) - f(x_2)| = |f(x_1) - a - (f(x_2) - a)| \le |f(x_1) - a| + |f(x_2) - a| < \frac{\epsilon}{3} + \frac{\epsilon}{3} = \frac{2 \epsilon}{3} \\
    |f(x_1) - f(x_2)| < \frac{2 \epsilon}{3} \\
    \sup_{x_1, x_2 \in B} |f(x_1) - f(x_2)| \le \frac{2 \epsilon}{3} < \epsilon \Rightarrow
    \omega(f, B) < \epsilon
\end{gather*}

\paragraph{2024.11.08 (3, 4)}
\subsubsection{Доказательство 1 $\Leftarrow$ 2}
\begin{gather*}
    \text{Критерий Коши} \Rightarrow \exists \lim \\
    \forall \epsilon > 0 ~ \exists B \in \mathcal{B}, \omega(f, B) < \epsilon \Rightarrow
    \exists \lim_{\mathcal{B}} f(x) = A, a \in \mathbb{R} \\
    \exists B_1 \in \mathcal{B}: \omega(f, B_1) < 1 \\
    \exists B_\frac{1}{2} \in \mathcal{B}: \omega(f, B_{\frac{1}{2}}) < \frac{1}{2} \\
    \dots \\
    \exists B_\frac{1}{n} \in \mathcal{B}: \omega(f, B_{\frac{1}{n}}) < \frac{1}{n} \\
    \text{Пусть } x_1 \in B_1 \\
    x_{\frac{1}{2}} \in B_{\frac{1}{2}} \\
    \dots \\
    x_{\frac{1}{n}} \in B_{\frac{1}{n}} \\
    \{x_n\} \rightarrow \{f(x_n)\} \\
    \text{Пусть } x \in B_{\frac{1}{n}} \\
    |f(x_n) - f(x_m) \le |f(x_n) - f(x)| + |f(x_m) + f(x)| < \frac{1}{n} + \frac{1}{m} \\
    f(x) \text{ является последовательностью, удовлетворяющей критерию Коши } \\
    \Rightarrow \text{ она сходится} \Rightarrow
    \exists \lim_{n \to \infty} f(x_n) = A \\
    \forall \epsilon > 0 ~ \exists B \in \mathcal ~ \forall x \in B: |f(x) - A| < \epsilon \\
\end{gather*}
"Зорич forever" (c) Родников


\subsection{Определение | Монотонные функции}
\[
    f: E \rightarrow \mathbb{R} \\
\]

\subsubsection{Определение | Возрастающая функция}
$f$ называется возврастающей на $E$, если
\[
\forall x_1, x_2 \in E: x_1 < x_2 \Rightarrow f(x_1) < f(x_2).
\]
Является строгомонотонной функцией.

\subsubsection{Определение | Неубывающая функция}
$f$ называтся неубывающей на $E$, если
\[
    \forall x_1, x_2 \in E: x_1 < x_2 \Rightarrow f(x_1) \le f(x_2).
\]

\subsubsection{Определение | Невозрастающая функция}
$f$ называтся невозрастающей, на $E$, если
\[
    \forall x_1, x_2 \in E: x_1 < x_2 \Rightarrow f(x_1) \ge f(x_2).
\]

\subsubsection{Определение | Убывающая функция}
$f$ называтся убывающей, на $E$, если
\[
    \forall x_1, x_2 \in E: x_1 < x_2 \Rightarrow f(x_1) > f(x_2).
\]
Является строкомонотонной функцией.


\section{Ограниченные функции}
\begin{gather*}
    f: E \rightarrow \mathbb{R} \\
    s = \sup E ~ (\text{если } E \text{ неограничено сверху, то пусть } s = +\infty) \\
    i = \inf E ~ (\text{если } E \text{ неограничено снизу,  то пусть } i = -\infty) \\
    i, s \text{ --- предельные точки множества } E \Rightarrow
    x \rightarrow s, x \in E; x \rightarrow i, x \in E
\end{gather*}

\subsection{Утверждение} 
\begin{gather*}
    \text{Пусть } f(x) \text{ --- неубывающая функция на } E, \\
    \text{тогда } \left( \exists \lim_{x \to s, ~ x \in E} f(x) = A, A \in \mathbb{R} \right) \Leftrightarrow
    \left( f(x) \text{ --- ограничена на } E \right)
\end{gather*}

\subsubsection{Доказательство $1 \Rightarrow 2$}
\begin{gather*}
    \exists \lim_{x \to s, ~ x \in E} f(x) = A \Rightarrow
    \forall \epsilon > 0 ~ \exists E_{\delta} = \{x, x \in E, s - \delta < s \le s\} \\
    \forall x \in E_{\delta}: |f(x) - A| < \epsilon \Rightarrow f(x) < A + \epsilon \\
    \forall x_1 \in E \setminus E_{\delta}: f(x_1) \le f(x) ~ (x \in E_{\delta}) \Rightarrow
    \forall x \in E: f(x) < A + \epsilon \rightarrow f(x) \le A \\
    f(x) \text{ --- ограничена на } E
\end{gather*}

\subsubsection{Доказательство $2 \Rightarrow 1$}
\begin{gather*}
    f(x) \text{ --- ограничена сверху на } E \\
    A_1 = \sup_{x \in E} f(x) \\
    \forall \epsilon > 0 ~ \exists x_0 \in E: A_1 - \epsilon < f(x_0) \le A_1 \\
    \forall x \in E, x > x_0: A_1 - \epsilon < f(x_0) < f(x) \le A_1 \\
    \forall \epsilon > 0 ~ \exists x_0 ~ \forall \underbrace{x \in E, x \in (x_0, s)}_{\text{элемент базы } x \rightarrow s, ~ x \in E}:
    A_1 - \epsilon < f(x) \le A_1 \Rightarrow
    |f(x) - A_1| < \epsilon \Rightarrow \exists \lim_{\underset{x \in E}{x \to s}} f(x) = A_1
\end{gather*}


\subsection{Утверждение}
\[
    \left( \exists \lim_{\underset{x \in E}{x \to i}} f(x) = A, A \in \mathbb{R} \right) \Leftrightarrow
    \left( f(x) \text{ ограничена снизу}\right)
\]


\subsection{Утверждение}
\[
    \left( \exists \lim_{\underset{x \in E}{x \to s}} f(x) = A, A \in \mathbb{R} \right) \Leftrightarrow
    \left( f(x) \text{ ограничена снизу}\right)
\]


\subsection{Утверждение}
\[
    \left( \exists \lim_{\underset{x \in E}{x \to i}} f(x) = A, A \in \mathbb{R} \right) \Leftrightarrow
    \left( f(x) \text{ ограничена сверху}\right)
\]


\subsubsection{Пример для всех Утверждений}
\begin{gather*}
    y = \arctg x ~~~ x \in \mathbb{R} ~~~ -\frac{\pi}{2} < y < \frac{\pi}{2} \\
    \exists \lim_{x \to +\infty} y =  \frac{\pi}{2} \\
    \exists \lim_{x \to -\infty} y = -\frac{\pi}{2}
\end{gather*}




\section{Предел композиции функции}
\subsection{Утверждения}
%\begin{gather*}
%\begin{enumerate}
%    \item f: Y \rightarrow \mathbb{R}, \mathcal{B}_Y, \exists \lim_{\mathcal{B}} f(y)
%    \item g: X \rightarrow Y, \mathcal{B}_X
%    \item \forall B_y \in \mathbb{B}_Y ~ \exists B_x \in \mathcal{B}_X: g(B_x) \subset B_y
%    (1) (2) (3) \Rightarrow \exists \lim_{\mathcal{B}_X} f(g(x)) = \lim_{\mathcal{B}_y} f(y)
%\end{enumerate}
%\end{gather*}

\subsubsection{Доказательство}
\begin{gather*}
    \text{Пусть } \lim_{\mathcal{B}_Y} f(y) = A \\
    \begin{rcases}
        \forall V(A) \exists B_y \in \mathcal{B}_y: f(B_y) \subset V(A) \\
        \exists B_x \in \mathcal{B}: g(B_x) \subset B_y
    \end{rcases} \Rightarrow
    \forall V(A) ~ \exists B_x \in \mathcal{B}_x: f(g(B_x)) \subset V(A) \Rightarrow 
    \exists \lim_{\mathcal{B}} f(g(x)) = A
\end{gather*}

\subsubsection{Пример}
\begin{gather*}
    \lim_{x \to 0} \frac{\sin 5x}{5x} = \lim_{y \to 0} \frac{\sin y}{y}, y = 5x \\
    \lim_{x \to 1} \frac{\sin (x - 1)}{x - 1} = 1
\end{gather*}

\subsection{...}
\begin{gather*}
    \lim_{x \to 0} (1 + x)^{\frac{1}{x}} = e \\
    \lim_{x \to \infty} (1 + \frac{1}{x})^x = e \\
    \lim_{x \to 0} \frac{\ln (1  + x)}{x} = 1 ~~~ \frac{\ln (1 + x)}{x} = \ln (1 + x)^{\frac{1}{x}} \\
    ~~~~~~~~~~~~~~~~~~~~~~~~~~~~~~~~~~~~~~ x \to 0
\end{gather*}


\section{Анализ бесконечно малых}

\subsection{Определение}
Функция

\subsection{Определение}
Функция $f(x)$ называется бесконечно малой по сравнению с $g(x)$ при $\mathcal{B}$, если
$f(x) = \alpha(x)g(x)$, где $\alpha(x)$ --- бесконечно малая при $\mathcal{B}$.
\begin{gather*}
    f(x) = \overline{\overline{o}} \text{ при } \mathcal{B} \\
    x^2 = \overline{\overline{o}}(x) ~~~ x \to 0 \\
    \tg x - \sin x = \overline{\overline{o}} ~~~ x \to 0 \\
    x = \overline{\overline{o}}(1) ~~~ x \to 0 \\
    1 - \cos x = \overline{\overline{o}}(x) ~~~ x \to 0
\end{gather*}

\paragraph{Примечание}
$\overline{\overline{o}}$ --- "о"  большая (из истории в аудио 2024.11.08 (4) 0:10:00-0:17:00)

\subsubsection{Пример}
\begin{gather*}
    \lim_{x \to 0} \frac{\sin x}{x} = 1 \\
    \frac{\sin x}{x} = 1 + \alpha(x) \\
    \alpha(x) \text{ --- бесконечно малая при } x \to 0 \\
    \sin x = x + \alpha(x)x \\
    \sin x = x(1 + \alpha(x))
\end{gather*}


\subsection{Определение}
\begin{gather*}
    f(x) = \underline{\underline{o}}(g(x)) \text{ при } \mathcal{B} \\
    \text{Если } f(x) = \varphi(x)g(x), \text{ где } \varphi \text{ --- финально ограничена при } \mathcal{B}
\end{gather*}


\subsection{Определение}
$f(x)$ и $g(x)$ называется величинами одного порядка при $\mathcal{B}$, если
$f(x) = O(g(x))$ и $g(x) = O(f(x))$


\subsection{Утверждение}
$f(x)$ и $g(x)$ являются величинами одного порядка при $\mathcal{B}$, если
\begin{gather*}
    \exists B\in \mathcal{B} ~~~ c_1, c_2, \text{ что } \forall x \in B: c_1|f(x)| \le |g(x)| \le c_2|f(x)| \\
    (c_1 > 0, c_2 > 0)
\end{gather*}

\subsubsection{Доказательство}
\begin{gather*}
    \text{Пусть } f(x), g(x) \text{ --- бесконечно малые при } \mathcal{B} \\
    c_1|f(x)| \le |g(x)| \le c_2|f(x)| \\
    c_1 \le \underbrace{\frac{|g(x)|}{|f(x)|}}_{\text{эта функция ограничена}} \le c_2 
    \Rightarrow \frac{1}{c_2} \le |\frac{f(x)}{g(x)}| \le \frac{1}{c_1} \\
    g(x) = O(f(x)); f(x) = O(g(x))
\end{gather*}

\paragraph{Замечание}
\begin{gather}
    \text{Если } \exists \lim_{\mathcal{B}} \frac{f(x)}{g(x)} = A, A \neq O,
    \text{то } f(x) \text{ и } g(x) \text{ --- величины одного порядка}
\end{gather}


\subsection{Определение}
\begin{gather*}
    f(x) \text{ и } g(x) \text{ --- бесконечно малые при } \mathcal{B} \\
    \text{Говорят, что } f(x) \sim g(x) (\sim \text{ --- Эквивалентность})
    \text{если } f(x) = \gamma(x) g(x), \lim_{\mathcal{B}} \gamma(x) = 1
\end{gather*}

\subsubsection{Замечание}
\begin{gather*}
    f(x) = (1 + \alpha(x)) \cdot g(x) \\
    \alpha(x) \text{ --- бесконечно малая при } \mathcal{B}
\end{gather*}

\subsubsection{Замечание}
...


\section{Эквивалентности}

\subsection{Основные эквивалентности}
\begin{enumerate}
    \item $\sin x \sim x$
    \item $\ln(1 + x) \sim x$
    \item $e^x - 1 \sim x$
    \item $(1 + x)^{\alpha} - 1 \sim \alpha x$
    \item $1 - \cos x \sim \frac{x^2}{2}$
\end{enumerate}

\subsubsection{Пример}
\begin{gather*}
    \lim_{x \to 0} \frac{1 - \cos x}{\frac{x^2}{2}} =
    \lim_{x \to 0} \frac{2 \sin^2 \frac{x}{2}}{\frac{x^2}{2}} =
    \lim_{x \to 0} \frac{\sin \frac{x}{2}}{\frac{x}{2}} \cdot \frac{\sin \frac{x}{2}}{...}
\end{gather*}


\subsection{Некоторые другие эквивалентности}
\begin{enumerate}
    \item $\tg x \sim x ~~~ x \to 0$
    \item $\sqrt{1 + x} - 1 \sim \frac{x}{2}$ 
\end{enumerate}


\subsection{Утверждение}
\begin{gather*}
    \lim_{\mathcal{B}} \frac{f(x)}{g(x)} = A \\
    f(x), g(x) \text{ --- бесконечно малые при } \mathcal{B} \\
    \text{Тогда если } f(x) \sim h(x) \text{ при } \mathcal{B}, \text{ то } \lim_{\mathcal{B}} \frac{h(x)}{g(x)} = A
\end{gather*}

\subsubsection{Доказательство}
\begin{gather*}
    f(x) \sim h(x) \Rightarrow f(x) = (1 + \gamma(x)) h(x),\gamma(x) \text{ --- бесконечно малая} \\
    \lim_{\mathcal{B}} \frac{h(x)}{g(x)} = ...
\end{gather*}



\section{Главная часть функции}

\subsection{Определение | Главная часть функции}
\begin{gather*}
    \text{Пусть } f(x) \text{ --- бесконечно малая при } x \to x_0 \\
    \text{Если } f(x) = C(x - x_0)^{\alpha} + h(x), \text{ где }
    h(x) = o (x -x_0)^{\alpha} \text{ при } x \to x_0, \\
    \text{то } C(x - x_0)^{\alpha} \text{ называется главной частью функции } f(x).
\end{gather*}

\subsubsection{Пример}
\begin{gather*}
    \lim_{x \to 0} \frac{\tg x - \sin x}{x^{\alpha}}
\end{gather*}



\paragraph{2024.11.15 (3)}
~\\
Преподаватель Байков Александр Евреньевич (доцент кафедры 802)

\subsection{Определение}
\begin{gather*}
    f(x) = o(g(x)), x \to a \\
    f(x) = \alpha(x) g(x), \\
    \text{где } \lim_{x \to 0} \alpha(x) = 0 \Leftarrow \lim_{x \to a} \frac{f(x)}{g(x)} = 0 \\
    ~\\
    f(x) = O(g(x)), x \to a \\
    f(x) = \beta(x) g(x), \\
    \text{где } \beta \text{ --- финально ограниченная функция при } x \to a
\end{gather*}


\section{Свойства $o$ $O$ (при $\mathcal{B}$)}

\subsection{Первое}
\begin{itemize}
    \item $o(f) \pm o(f) = o(f)$
    \item $\alpha_1 f \pm \alpha_2 f = (\alpha_1 \pm \alpha_2) f$
    \item $o(f) \pm O(f) = O(f)$
    \item $O(f) \pm O(f) = O(f)$
\end{itemize}

\paragraph{Проколотая окрестность $a$}
\[
    \overset{o}{U}_{\delta}(a) = U_{\delta}(a) \setminus \{a\}
\]
\includegraphics{images/2024-11-15/02.png}


\subsection{Второе}
\begin{itemize}
    \item $o(o(f)) = o(f)$
    \item $o(O(f)) = o(f)$
    \item $O(o(f)) = o(f)$
    \item $O(O(f)) = O(f)$
\end{itemize}

\begin{gather*}
    \underbrace{\beta(x)}_{\text{огр.}} \underbrace{\alpha(x)}_{\text{б. м.}} f(x) = \gamma(x) f(x) \\
    \gamma(x) = \beta(x) \alpha(x) \text{ --- бесконечно малая}
\end{gather*}


\subsection{Третье}
\begin{itemize}
    \item $o(f + o(f)) = o(f)$
    \item $o(f + O(f)) = o(f)$
\end{itemize}


\subsection{Четвёртное}
\begin{itemize}
    \item $o(f) o(g) = o(fg)$
    \item $O(f) o(g) = o(fg)$
    \item $f o(g) = o(fg)$
\end{itemize}

\subsubsection{Пример}
\begin{itemize}
    \item $o(1)$ --- бесконечно малая функция
    \item $O(1)$ --- финально ограниченная функция
\end{itemize}

\begin{gather*}
    f(x) = g(x) + o(g(x)), x \to a \\
    g \cancel{\equiv} 0 \text{ в проколотой окрестности } a \\
    g \text{ --- главная часть функции } f \text{ при } x \to a \\
    \text{Обычно } g(x) = C(x - a)^{\alpha}, C \neq 0, \alpha \in \mathbb{R} \\
    g \text{ --- степенная главная часть} \Leftrightarrow \\
    \Leftrightarrow f(x) \sim g(x), 
\end{gather*}


\section{Фото 8 | Пример}
\paragraph{Дано}
\[
    f(x) = \tg x - \sin x, x \to 0 \\
\]

\paragraph{Найти степенную главную часть}
\[
    f(x) = cx^{\alpha} + o(x^{\alpha})
\]

\paragraph{Не решение}
\begin{gather*}
    f(x) = x + o(x) - (x + o(x)) = o(x) - o(x) = o(x), x \to 0 \\
    f(x) = \tg x (1 - \cos x) = (x + o(x)) \left( \frac{x^2}{2} + o(x^2) \right) = \\
    = \frac{1}{2} x^3 + x o(x^2) + \frac{x^2}{2} o(x) + o(x) o(x^2) =
    \frac{1}{2} x^3 + o(x^3) + o(x^3) + o(x^3) = \\
    = \frac{1}{2} x^3 + o(x^3), x \to 0 \Rightarrow
    g(x) = \frac{1}{2} x^3 \text{ --- главная часть}
\end{gather*}



\section{Непрерывность}
\includegraphics[scale=0.5]{images/2024-11-15/01.png}

\subsection{Определение | Непрерывность функции в точке по Коши}
\[
    f: E \to \mathbb{R}, a \in E
\]

$f$ непрерывна в точке $a$, если:
\begin{gather*}
    \forall \epsilon > 0 ~ \exists \delta = \delta(\epsilon) > 0 ~ \forall x \in E \\
    |x - a| < \delta \Rightarrow |f(x) - f(a)| < \epsilon
\end{gather*}
В определении предела $|x - a| < \delta \Rightarrow |f(x) - A| < \epsilon$


\subsection{Определение}
\begin{gather*}
    \forall V(f(a)) ~ \exists U(a) ~ \underbrace{f(U_E(a)) \subset V(f(a))}_{\forall x \in U_E(a) ~~~ f(x) \in V(f(a))} \\
    U_E(a) := U(a) \cap E \\
    f(U_E(a)) \text{ --- образ}
\end{gather*}
\includegraphics[scale=0.3]{images/2024-11-15/03.png}




\subsection{Определение | Непрерывность функции в точке по Гейне}
\[
    f: E \rightarrow \mathbb{R}, a \in E
\]

\subsubsection{Случай 1}
$f$ непрерывна в точке $a$, если:
\[
    \forall \{x_n\}^{\infty}_{n = 1} \supset E, \lim_{n \to \infty} x_n = a \Rightarrow \lim_{n \to \infty} f(x_n) = f(a)
\]

\subsubsection{Случай 2}
$a$ --- предельная точка множества $E$:
\begin{gather*}
    \forall \delta > 0 ~~ \overset{o}{U}_{\delta}(a) \cap E \neq \varnothing ~\Leftrightarrow~
    \forall \delta > 0 ~~ U_{\delta}(a) \cap E \text{ --- бесконечное множество }
\end{gather*}

\subsubsection{Пример}
\begin{gather*}
    f: E \to \mathbb{R} \\
    f(x) \equiv c \Leftrightarrow \forall x \in E f(x) = c.
\end{gather*}

\subsubsection{Случай 3}
$a$ --- внутренняя точка множества $E$:
\begin{gather*}
    \exists \delta_0 > 0 ~ U_{\delta} \subset E.
\end{gather*}

\subsubsection{Связь непрерывности и предела}
Если $f$ непрерывна в точке $a$ --- предел в точке для $E$, то
\[
    \exists \lim_{x \to a} f(x) = f(a), x \to a, x \in E.
\]

% \subsection{Замечание}
% Если $f$ не является непрерывной в точке $a$, то она разрывна в точке $a$. 

% \subsubsection{Случай 1}
% $a$ --- изолированная точка множества $E$ \\

% \begin{gather*}
%     \exists \delta_0 > 0 ~ U_{\delta_0} \cap E = (a - \delta_0, a + \delta_0) = \{a\}
% \end{gather*}
% (график 3 (прямая с буквой Е над полосками на прямой))
% \begin{gather*}
%     f: E \rightarrow \mathcal{R} \\
% \end{gather*}



\section{Функция Дирихле (фото 14, 15)}
\begin{gather*}
    D(x) =
    \begin{cases}
        1, x \in \mathbb{Q} \\
        0, x \in \mathbb{R} \setminus \mathbb{Q}
    \end{cases} \\
    D: \mathbb{R} \rightarrow \mathbb{R} \\
    a \in \mathbb{R} \setminus \mathbb{Q} \\
    x_n = a - \frac{1}{n} \underset{n \to \infty}{\longrightarrow} a \\
    \{y_n\}_{n = 1} \subset \mathbb{Q}, y_n \underset{n \to a}{\longrightarrow} \\
    \lim_{n \to \infty} D(x_n) = \lim_{n \to \infty} 0 = 0 \\
    \lim_{n \to \infty} D(y_n) = \lim_{n \to \infty} 1 = 1 \\
    \nexists \lim_{x \to n} D(x) \Rightarrow D \text{ разрывна в точке } a
\end{gather*}



\section{фото 16, 17}

\begin{gather*}
    x_n^{'}  \underset{n \to \infty}{\longrightarrow} a \in \mathbb{R} \setminus \mathbb{Q} \\
    x_n^{''} \underset{n \to \infty}{\longrightarrow} a \\
    \lim_{n \to \infty} R(x_n) = 0 \Leftarrow
    \begin{cases}
        \underset{n \to \infty}{\lim} ~ R(x_n^{''}) = \underset{n \to \infty}{\lim} ~ 0 = 0 \\
        \underset{n \to \infty}{\lim} ~ R(x_n^{'})  = \underset{n \to \infty}{\lim} ~ R \left( \frac{k_n}{m_n} \right) =
        \underset{n \to \infty}{\lim} ~ \frac{1}{m_n} = 0
    \end{cases} \\
    m_n \to \infty, n \to \infty, \text{ так как } x_n^{'} \to a \in \mathbb{R} \setminus \mathbb{Q}, n \to \infty
\end{gather*}



\section{Классификация точек разрыва}

$a$ --- внутренняя точка $E$, $f: E \rightarrow \mathbb{R}$

\subsection{Определение}
Функция $f$ имеет устранимый разрыв в точке $a$, если
\[
    \exists \lim_{x \to a} f(x) \neq f(a)
\]

\paragraph{Переопределение}
\begin{gather*}
    \overset{\sim}{f} (x) =
    \begin{cases}
        f(x), x \in E \setminus \{a\} \\
        \underset{t \to a}{\lim} ~ f(t), x = a
    \end{cases}
\end{gather*}


\subsection{Определение}
Функция $f$ имеет разрыв первого рода в точке $a$, если
\begin{gather*}
    \exists \lim_{x \to a + 0} f(x) = f_+(a) \\
    \exists \lim_{x \to a - 0} f(x) = f_-(a) \\
    f_+(a) \neq f_-(a)
\end{gather*}

\subsubsection{Пример}
\begin{gather*}
    f(x) = sign(x) \\
    x = 0 \text{ --- точка разрыва первого порядка}
\end{gather*}


\paragraph{2024.11.22 (3, 4)}
\subsection{Определение}
Разрывом второго рода называется разрыв, который не является разрывом первого рода.

\subsubsection{Пример}
\[
    y = \frac{1}{x}
\]
\begin{gather*}
    \lim_{x \to 0-} \frac{1}{x} = -\infty \\
    \lim_{x \to 0+} \frac{1}{x} = +\infty \\
    x = 0 \text{ --- точка разрыва второго рода}
\end{gather*}
\includegraphics[scale=0.1]{images/2024-11-22/01.png}

\subsubsection{Пример}
\[
    y = e^{\frac{1}{x - 1}}
\]
\begin{gather*}
    \lim_{x \to 1-} e^{\frac{1}{x - 1}} = 0 \\
    \lim_{x \to 1+} e^{\frac{1}{x - 1}} = +\infty \\
    \text{разрыв второго рода}
\end{gather*}
\includegraphics[scale=0.5]{images/2024-11-22/02.png}

\subsubsection{Пример}
\[
    y = \sin \frac{1}{x}
\]
\begin{gather*}
    \nexists \lim_{x \to 0+} \sin \frac{1}{x} \\
    \nexists \lim_{x \to 0-} \sin \frac{1}{x} \\
    \text{разрыв второго рода}
\end{gather*}
\includegraphics[scale=0.3]{images/2024-11-22/03.png}


\subsection{Определение | Точка устранимого разрыва}
\begin{gather*}
    \text{Пусть } f: E \rightarrow \mathbb{R}, a \in E, f \notin C(a) \\
    \text{Если } \exists \overset{\sim}{f}: E \rightarrow \mathbb{R}, 
    \overset{\sim}{f} \in E \text{ при } x \neq a, \overset{\sim}{f}(x) = f(x), \\
    \text{то } x = a \text{ называется точкой устранимого разрыва}
\end{gather*}

\subsubsection{Пример}
\begin{gather*}
    y = | \underbrace{signum}_{= f(x)} ~ x | = 
    \begin{sqcases}
        1, x \neq 0 \\ 
        0, x = 0
    \end{sqcases}
\end{gather*}

\subsubsection{Пример}
\begin{gather*}
    \text{Пусть } a \text{ --- точка устранимого разрыва} \\
                \exists \lim_{x \to a}  \overset{\sim}{f}(x) = \overset{\sim}{f}(a) \Rightarrow \\
    \Rightarrow \exists \lim_{x \to a-} \overset{\sim}{f}(x) = \overset{\sim}{f}(a),
                \exists \lim_{x \to a+} \overset{\sim}{f}(x) = \overset{\sim}{f}(a) \\
\end{gather*}
    
\subsubsection{Следствие}
Устранимый разрыв --- это разрыв первого рода, причём пределы фцнкии слева и справа равны между собой.



\section{Функция Римана}

\subsection{Определение}
\begin{gather*}
    R(x) =
    \begin{cases}
        \frac{1}{n}, x = \frac{m}{n} \text{ --- несократимая дробь} \\
                  0, x \in \mathbb{R} \setminus \mathbb{Q}
    \end{cases}
\end{gather*}

\subsubsection{Пример}
\paragraph{Дано}
\[
    \text{Если } x_0 \in \mathbb{R} \setminus \mathbb{Q}, \text{ то } R(x) \in C(x_0)
\]
\paragraph{Доказательство}
\begin{gather*}
    \epsilon > 0 ~~~ M = \{x, R(x) \ge \epsilon\} \Rightarrow \text{если } x \in M, \text{ то } x \in \mathbb{Q} \\
    M = \left\{ \frac{m}{n}; \underbrace{\frac{1}{n} \ge \epsilon} \right\} \\
    ~~~~~~~~~~~~~ n \le \frac{1}{\epsilon} \\
    \text{В любом ограниченном промежутке находится конечное число элементов множества } M
\end{gather*}
\begin{gather*}
    \text{Пусть } x_0 \in \mathbb{R} \setminus \mathbb{Q} \\
    \overset{\sim}{x} \text{ --- наиболее близкий к } x_0 \text{ элемент } M \\
    \delta = \frac{\left| x_0 - \overset{\sim}{x} \right|}{2} \\
    \text{в промежутке } (x_0 - \delta; x_0 + \delta) \text{ нет элементов } M, \\
    \text{то есть } \forall x \in (x_0 - \delta; x_o + \delta): R(x) < \epsilon \\
    \exists \lim_{x \to x_0} R(x) = R(x_0) = 0
\end{gather*}
\includegraphics[scale=0.3]{images/2024-11-22/04.png}


\section{Локальные свойства непрерывных функций}

\subsection{Утверждение}
\begin{gather*}
    \text{Пусть } f: E \rightarrow \mathbb{R}, a \in E, \underbrace{f(x) \in C(a)}_{\text{функция непрерывна в точке } a} \\
    \text{тогда } \exists U_E(a), \text{ в которой } f(x) \text{ --- ограничена}
    \left( U_E(a) = \underbrace{U(a)}_{\text{окрестность точки } a} \cap E \right)
\end{gather*}

\subsubsection{Доказательство}
\begin{gather*}
    \forall \epsilon > 0 ~ \exists U_E(a), \forall x \in U_E(a): |f(x) - f(a)| < \epsilon \Rightarrow \\
    \Rightarrow \underbrace{f(a) - \epsilon}_{\text{миноранта}} < f(x) < \underbrace{f(a) + \epsilon}_{\text{мажоранта}}
\end{gather*}


\subsection{Утверждение}
\begin{gather*}
    \text{Пусть } f: E \rightarrow \mathbb{R}, a \in E, f \in C(a), f(a) > 0 \\
    \text{тогда } \exists U_E(a), \forall x \in U_E(a): f(x) > 0
\end{gather*}

\subsubsection{Доказательство}
\begin{gather*}
    \text{Пусть } \epsilon = \frac{f(a)}{2} \\
    \text{тогда } \exists U_E(a), \forall x \in U_E(a): \left| f(x) - f(a) \right| < \epsilon \Rightarrow \\
    \Rightarrow f(a) - \epsilon < f(x) \Rightarrow \\
    \Rightarrow 0 < \frac{f(a)}{2}< f(x)
\end{gather*}


\subsection{Утверждение | Арифметические свойства непрерывных функций}
\begin{gather*}
    \begin{cases}
        f(x) \in C(a) \\
        g(x) \in C(a)
    \end{cases} \Rightarrow
    \begin{cases}
        f(x) + g(x) \in C(a) \\
        f(x) \cdot g(x) \in C(a) \\
        \frac{f(x)}{g(x)} \in C(a)$, если $g(a) \neq 0
    \end{cases}
\end{gather*}

\subsubsection{Доказательство}
Доказательства каждого свойства можно провести по соответствующим свойствам пределов.


\subsection{Утверждение}
\begin{gather*}
    \begin{cases}
        f: E \rightarrow \mathbb{R}, f \in C(a), a \in E \\
        g: U \rightarrow E,          g \in C(a), g(b) = a
    \end{cases} \Rightarrow
    h(x) = f(g(x)) \in C(b)
\end{gather*}
\includegraphics[scale=0.5]{images/2024-11-22/05.png}

\subsubsection{Доказательство}
Для доказательства достаточно сослаться на теорему о пределе композиций двух функций.



\section{Глобальные свойства непрерывных функций}

\subsection{Теорема (Больцано-Коши)}
\begin{align*}
    \begin{rcases}
        f(x) \in C([a, b]) \\
        f(a) \cdot f(b) < 0
    \end{rcases} \Rightarrow
    \exists c \in [a, b], f(c) = 0
\end{align*}
\includegraphics[scale=1]{images/2024-11-22/08.png}

\subsubsection{Доказательство}
Поделим отрезок $[a, b]$ пополам:
\begin{enumerate}
    \item $f(\frac{a + b}{2}) = 0$, то $c = \frac{a + b}{2}$
    \item $f(\frac{a + b}{2}) \neq 0$, то значение функции $f(x)$ на концах одной из половинок имеют разные знаки. \\
    Обозначим эту половинку как $[a_1, b_1]$. \\
    Для отрезка $[a_1; b_1]$ повторим процедуру. \\ 
    В результате либо найдм точку с, либо найдём отрезок $[a_2, b_2]$, обладающий тем же свойством. \\
    Будем повторять эту процедуру либо до нахождения точки $c$, либо либо до бесконечности. \\
    Предположим, пришлось выполнять процедуру до бесконечности. \\
    В результате получаем последовательность вложенных отрезков. \\
\end{enumerate}
\begin{gather*}
    [a_2, b_2], [a_3, b_3], \dots, [a_n, b_n], \dots \\
    \lim_{n \to \infty} (b_n - a_n) = 0 \\
    \exists! c \in [a_n, b_n], n \in \mathbb{N} \\
    \exists \lim_{n \to \infty} a_n = c ~~~ f(a_n) < 0 \\
    \exists \lim_{n \to \infty} b_n = c ~~~ f(b_n) > 0 \\
    f(x) \in C(c) ~~~~ \exists \lim_{x \to c} f(x) = f(c) \\
    \lim_{n \to \infty} f(a_n) = f(c) \le 0, \lim_{n \to \infty} f(b_n) = f(c) \ge 0
\end{gather*}

\subsubsection{Следствие}
\begin{gather*}
    \text{Пусть } f(x): [a, b] \Rightarrow \mathbb{R}, ~~ 
    f(x) \in C([a, b]), ~~ 
    a_1, b_1 \in [a, b], ~~ 
    f(a_1) = A, f(b_1) = B, A < B, \\
    \text{тогда } \forall C \in [A, B] ~ \exists c \in [a_1, b_1], f(c) = C 
\end{gather*}
\includegraphics[scale=0.5]{images/2024-11-22/07.png}

\subsubsection{Доказательство}
\begin{gather*}
    F(x) = f(x) - C ~~~ F(a_1) < 0 ~~~ F(b_1) > 0 \rightarrow \\
    \rightarrow \text{Больцано-Коши: } \exists c, F(c) = 0 \Rightarrow f(c) = C    
\end{gather*}

\subsubsection{Примеры нарушения теоремы}
\dots


\subsection{Теорема | Вейерштрасс о максимальном значении}

Пусть $f(x) \in C([a, b])$, тогда
\begin{enumerate}
    \item $f(x)$ --- ограничена на $[a, b]$
    \item Если $M = \sup_{[a, b]} f(x)$, $m = \inf_{[a, b]} f(x)$, \\
    то $\exists c, C \in [a, b], f(c) = m, f(C) = M$
\end{enumerate}

\subsubsection{Доказательство}
\begin{gather*}
    \epsilon > 0 ~~~ \forall x_o \in [a, b] ~ \exists U_{[a, b]}(x_0) ~ \forall x \in U_{[a, b]}: |f(x) - f(x_0)| < \epsilon \\
    \{U(x_0), x_0 \in [a, b]\} \text{ --- покрытие отрезка} \\
    \text{Выделим из этого покрытия конечное подпокрытие: } \{U(x_{01}), U(x_{02}), \dots, U(x_{0N})\} \\
    \forall x \in U(x_{0k}): |f(x) - f(x_{0k})| < \epsilon \Rightarrow f(x) < f(x_{0k}) + \epsilon \\
    \text{Пусть } Z = \max \{f(x_{0k}, k = 1, 2, \dots, N)\}, \text{тогда } \forall x \in [a, b]: f(x) < Z + \epsilon
\end{gather*}
Тем самым мы доказали, что функция $f(x)$ ограничена сверху. Аналогично доказывается, что функция ограничивается снизу. \\
\includegraphics[scale=0.5]{images/2024-11-22/09.png}


\subsubsection{Доказательство}
\begin{gather*}
    \text{Пусть } M = \sup_{[a, b]} f(x), ~~~ \forall x \in [a, b]: f(x) < M, \\
    \text{тогда } \forall \epsilon > 0 ~ \exists y \in [a, b]: f(y) > M - \epsilon \\
    F(x) = \frac{1}{M - f(x)}, ~~~ F(x) \in C([a, b]) \\
    F(y) = \frac{1}{M - f(y)} > \frac{1}{\epsilon} \Rightarrow F(y) \text{ неограничена --- absurdum!} \Rightarrow
    \exists c, f(c) = M 
\end{gather*}
Аналогично доказывается, что наименьшее значение тоже достигается.


\subsubsection{Следствие} 
Область значений непрерывной или на отрезке функции есть отрезок.

\subsubsection{Примеры}
\includegraphics[scale=0.3]{images/2024-11-22/10.png}



\section{Равномерная непрерывность}
Пусть $f: E \rightarrow \mathbb{R}$

\subsection{Определение}
Функция $f(x)$ называется равномерно непрерывной на $E$, если
\[
    \forall \epsilon > 0 ~ \exists \delta(\epsilon) > 0 ~ \forall x_1, x_2 \in E, |x_1 - x_2| < \delta \Rightarrow |f(x_1) - f(x_2)| < \epsilon
\]


\subsection{Утверждение | Связь определения с обычной непрерывностью}
Если $f$ равномерно непрерывна на $E$, то $f$ непрерывна в любой точке $E$.

\subsubsection{Доказательство}
\begin{gather*}
    \text{Пусть } x_o \in E ~ \forall \epsilon > 0 ~ \exists \delta(\epsilon) > 0 ~ \forall x \in E,
    |x - x_0| < \delta: |f(x) - f(x_0)| < \epsilon, \\
    \text{ то есть } \exists \lim_{x \to x_0} f(x) = f(x_0), f(x) \in C(x_0) 
\end{gather*}
Из непрерывности не следует равномерная непрерывность.

\subsubsection{Пример}
\begin{gather*}
    y = x^2, x \in \mathbb{R} \\
    y = x^2 \in C(\mathbb{R}) \\
    x_1 = \sqrt{n}, ~ x_2 = \sqrt{n + 1} \\
    f(x_2) - f(x_1) = 1 \\
    x_2 - x_1 = \sqrt{n + 1} - \sqrt{n} \\
    \lim_{n \to \infty} (\sqrt{n + 1} - \sqrt{n}) = 
    \lim_{x \to \infty} \frac{1}{\sqrt{n + 1} + \sqrt{n}} \\
    y = x^2 \text{ не является равномерно непрерывной на } \mathbb{R}
\end{gather*}

\subsubsection{Пример}
\begin{gather*}
    y = \frac{1}{x}, x \in (0; 1) \\
    x_1 = \frac{1}{n + 1} ~~ y_1 = n + 1 \\
    x_2 = \frac{1}{n}     ~~ y_2 = n \\
    y_1 - y_2 = 1 \\
    \lim_{n \to \infty} (x_2 - x_1) = 
    \lim_{n \to \infty} (\frac{1}{n} - \frac{1}{n + 1}) =
    \lim_{n \to \infty} \frac{1}{n(n + 1)} = 0 \\
    y = \frac{1}{x} \text{ не является равномерно непрерывной на } (0; 1)
\end{gather*}



\subsection{Теорема равномерной непрерывности}
Если функция непрерывна на отрезка, то она равномерно непрерывна на этом отрезке.


\end{document}